% Fate's Edge -- Complete Grimoire
% A Thematic Guide for Cantors, Invokers, Runekeepers, and Summoners
% Compile with pdflatex

\documentclass[10pt,a5paper]{book}
\usepackage[top=1.5cm,bottom=1.5cm,left=1.2cm,right=1.2cm]{geometry}
\usepackage[T1]{fontenc}
\usepackage{lmodern}
\usepackage{titlesec}
\usepackage{tcolorbox}
\tcbuselibrary{skins,breakable}
\usepackage{multicol}
\usepackage{hyperref}
\hypersetup{colorlinks=true,linkcolor=blue,pdfencoding=auto}
\usepackage{array}
\usepackage{longtable}
\usepackage{epigraph}
\setlength{\epigraphwidth}{0.7\textwidth}
\renewcommand{\epigraphflush}{flushright}
\usepackage{lettrine}

% Chapter styling
\titleformat{\chapter}[display]
{\Huge\bfseries\scshape}
{\flushright\rule{0.2\textwidth}{0.5pt}\\\vspace{-1cm}}
{0pt}
{\hfill}
\titlespacing*{\chapter}{0pt}{-20pt}{40pt}

% Section styling
\titleformat{\section}
{\Large\bfseries}
{\thesection}{0.5em}{}
\titlespacing*{\section}{0pt}{12pt}{6pt}

% Subsection
\titleformat{\subsection}
{\normalsize\bfseries\itshape}
{\thesubsection}{0.5em}{}
\titlespacing*{\subsection}{0pt}{8pt}{4pt}

% Custom environments
\newtcolorbox{patronbox}[2][]{
    colback=gray!5,
    colframe=gray!70,
    fonttitle=\bfseries,
    title=⧫ #2 ⧫,
    enhanced,
    attach boxed title to top left={yshift=-2mm,xshift=3mm},
    boxed title style={colback=gray!20},
    #1
}

\newtcolorbox{ritualnote}{
    colback=yellow!5,
    colframe=orange!70,
    before skip=0.5cm,
    after skip=0.5cm
}

\newtcolorbox{warningbox}{
    colback=red!5,
    colframe=red!70,
    before skip=0.5cm,
    after skip=0.5cm
}

\newtcolorbox{spellcard}[2][]{
    breakable,
    colback=gray!5,
    colframe=black!50,
    fonttitle=\bfseries\large,
    title=#2,
    enhanced,
    attach boxed title to top left={yshift=-3mm,xshift=3mm},
    boxed title style={colback=white},
    #1
}

\newtcolorbox{spiritcard}[2][]{
    breakable,
    colback=blue!5,
    colframe=blue!70,
    fonttitle=\bfseries\large,
    title=#2,
    enhanced,
    attach boxed title to top left={yshift=-3mm,xshift=3mm},
    boxed title style={colback=blue!20},
    #1
}

% Tags for magic
\newcommand{\ritualtag}[1]{\textsc{#1}}
\newcommand{\dangerous}[1]{\textcolor{red}{\textbf{#1}}}
\newcommand{\tag}[1]{\textsc{#1}}

% Spell command
\newcommand{\spell}[5]{ % name, tags, dv, effect, backlash
    \begin{spellcard}{#1}
    \textbf{TAGS:} #2 \\
    \textbf{DV:} #3 \\
    \textbf{Effect:} #4 \\
    \textbf{Backlash (1+ SB):} #5
    \end{spellcard}
    \vspace{0.5cm}
}

\begin{document}

\frontmatter
\thispagestyle{empty}
\vspace*{2cm}
\begin{center}
{\Huge\scshape The Gray Wanderer's Grimoire}\\[0.5cm]
{\Large A Complete Guide to Cantors, Invokers, Runekeepers, and Summoners}\\[1cm]
\rule{0.6\textwidth}{0.5pt}\\[1cm]
{\itshape "You want power? Fine. But don"t come crying when the debt comes due."}
\end{center}
\vfill
\clearpage

\tableofcontents
\clearpage

\mainmatter

\chapter{Introduction -- A Voice from the Crossroads}
\epigraph{"I have worn the chains of a dozen Patrons and broken most of them. Twice."}{\textit{--- The Gray Wanderer}}

This grimoire is not a rulebook. It is a \textbf{companion} -- a set of lenses through which to see the magical paths of \textit{Fate's Edge}. We will explore:

\begin{itemize}
    \item The \textbf{lore and roleplay} of Cantors, Invokers, Runekeepers, and Summoners -- who they are, how they think, why they risk annihilation.
    \item How to handle \textbf{Obligation} -- not as a punishment, but as a story engine. You will find patron‑specific ideas, intrusions, and acts of service.
    \item \textbf{Rituals without a focus} -- when your Codex is ash, your Symbol cracked, your voice alone. Components matter again.
    \item The art of \textbf{Summoning} -- binding spirits to will, from the deep stone to the salt foam.
\end{itemize}

The voice you hear belongs to no one and everyone -- a wanderer who has stood at the crossroads of every Patron's attention. Call me the \textbf{Gray Wanderer}. These notes are my legacy. Use them, lose them, or burn them for warmth. I won"t mind.

\clearpage

\chapter{The Weight of Names -- How the Peoples of Fate's Edge See Magic}

\epigraph{"I have carried my power through a dozen cultures and been called a dozen things: saint, sinner, witch, warlock, fraud, fool. Every people has a word for what I do. None of them are quite right. None of them are quite wrong."}{\textit{--- The Gray Wanderer}}

Magic in Fate's Edge is not a single practice. It is a family of arguments with reality -- and every culture argues differently. Some embrace the structured debt of Patrons. Others whisper to spirits at crossroads. A few have learned to bend the world with nothing but a disciplined mind. What follows is not a definitive guide -- I have been corrected by village grandmothers and high priests alike -- but a wanderer's impressions, gathered over too many years on too many roads.

\section{Aeler -- The Deep Accountants}

\lettrine{T}{he} dwarves of Aeler do not speak of "magic." They speak of "load," "breath," and "ledger." A Rite is not a miracle -- it is a contract with stone, sealed by proper measure and witnessed by a bell. A Runekeeper is an engineer of the invisible, and a witch is a mender of cracks in the world's infrastructure.

\subsection*{What They Respect}
\begin{itemize}
    \item \textbf{Runekeeper:} The ideal Aeler mage. A Codex carved in slate, a Thiasos that chitters in the dark, and a precise understanding of Obligation. The mountain respects a debt you can count.
    \item \textbf{Witch (Hearth craft):} The women who keep the cisterns sweet and the gas vents predictable are not feared -- they are salaried. Aeler employs witches as "Breath Wardens" and "Cistern Mothers."
\end{itemize}

\subsection*{What They Mistrust}
\begin{itemize}
    \item \textbf{Cantor:} Too public. Too emotional. The deep stone does not like loud noises it cannot measure.
    \item \textbf{Free Caster:} An abomination. Power without a ledger is a cave-in waiting to happen.
    \item \textbf{Psion:} Rare, but when one appears, the Aeler assign them to the Deep Vaults -- where their mind-reading is useful and their instability is contained.
\end{itemize}

\subsection*{The Breath-Bond}
Every Aeler mage, regardless of path, must learn the \textbf{Breath Bond}: a meditation that converts psychic strain into measured breaths. A Psion in Aeler carries a bell on a leather cord. Each Mental Strain point is one bell-ring. When the bell rings three times too fast, a Corrector appears.

\subsection*{Famous Aeler Magic}
The \textbf{Spine Underway} -- a tunnel that remembers every traveler who has passed through it. Runekeepers maintain it by carving new Rites into the walls each century. Witches keep the air sweet. And the stone itself? It is said to be the oldest Summoner, awake and patient.

\section{Aelaerem -- The Hearth-and-Hollow}

\lettrine{T}{he} halflings of Aelaerem do not build temples to Patrons. They leave butter on cup-mark stones and hang red thread on hawthorn. Their magic is domestic -- and therefore invisible to those who look only for grand gestures.

\subsection*{What They Respect}
\begin{itemize}
    \item \textbf{Witch:} The Aelaerem are the heartland of witchcraft. Every village has a Hedge-Mother who knows which knot to tie for a safe birth and which bell to ring to keep the Hollow from swallowing a child.
    \item \textbf{Cantor:} Folk songs are the oldest grimoire. A Cantor who sings the old tunes and respects the "no ninth" rule is welcome at any hearth.
    \item \textbf{Summoner:} But only brownies, field-wights, and the small dead. Aelaerem summoners do not call demons -- they ask the ancestors to move the firewood.
\end{itemize}

\subsection*{What They Mistrust}
\begin{itemize}
    \item \textbf{Invoker:} Carrying Symbols from multiple Patrons is seen as greedy. The Hollow does not forgive divided loyalty.
    \item \textbf{Psion:} "Reading thoughts without asking" is a violation of hospitality. A known psion will find doors closing and cup-mark stones left dry.
\end{itemize}

\subsection*{The Hollow Tithe}
Once a year, the Aelaerem pay the \textbf{Hollow Tithe}: a basket of the year's best bread, left at the edge of a barrow. Witches conduct the rite. No one speaks during the walk back. I have done it three times. Something was always following. Something always stopped at the village gate.

\section{Aelinnel -- The Bright Weird}

\lettrine{T}{he} gnomes of Aelinnel live in a world where numbers have manners and doors sometimes lead to last Tuesday. Their magic is not chaotic -- it is \textit{algebraic}. They have names for equations that other cultures do not know exist.

\subsection*{What They Respect}
\begin{itemize}
    \item \textbf{Free Caster:} The gnomes invented the TAGS system, though they call it "Contextual Recursion." A free caster who can improvise a spell on the spot is a philosopher, not a danger.
    \item \textbf{Psion:} The mind is the ultimate context. Aelinnel psions are trained from childhood to "count the thoughts" and "measure the echo."
    \item \textbf{Summoner:} Spirits are just data from a previous iteration of reality. Treat them politely and document everything.
\end{itemize}

\subsection*{What They Mistrust}
\begin{itemize}
    \item \textbf{Runekeeper:} Too rigid. A Rite carved in stone cannot adapt to a changing context. The gnomes prefer the flexibility of witch-knots and psionic whispers.
    \item \textbf{Cantor:} A song that cannot be written in two ledgers (Said and Meant) is dangerous. The gnomes have banned three hymns in the last decade.
\end{itemize}

\subsection*{The Green Gate}
The most famous Aelinnel threshold -- a sea-arch that opens at certain tides to a forest not on any map. Witches and Free Casters use it to travel between realities. Psions use it to test their mental resilience. The gnomes charge a toll: one truth you have never told anyone, and one lie you have never been caught telling.

\section{Linn -- The Storm-Tested}

\lettrine{T}{he} northmen of Linn do not ask permission from Patrons. They make oaths -- and the sea witnesses. Their magic is carved into keels, braided into ropes, and whispered into the ears of drowning men.

\subsection*{What They Respect}
\begin{itemize}
    \item \textbf{Summoner (Selkie and Drowned):} A Linn tide-witch who can call a selkie for navigation or a drowned captain for warning is worth her weight in whalebone.
    \item \textbf{Cantor:} Sea shanties are weather magic. A Cantor who sings the right note at the right moment can calm a storm -- or call one.
    \item \textbf{Witch (Wake-Judge):} The Linn have formalized witchcraft into a legal system. Wake-Judges arbitrate disputes at the water's edge, using promise clocks to bind oaths.
\end{itemize}

\subsection*{What They Mistrust}
\begin{itemize}
    \item \textbf{Invoker:} Borrowing power from a Patron is seen as weakness. A real northman makes his own luck.
    \item \textbf{Psion:} Too quiet. The Linn prefer their magic loud enough to hear over the wind.
\end{itemize}

\subsection*{The Thing-Bell}
Every Linn longship has a bell forged from the melted-down sword of the ship's first captain. When a Runekeeper or Witch needs to seal an oath, they ring the Thing-Bell. The sound carries across water. The dead listen.

\section{Ykrul -- The Weather Readers}

\lettrine{T}{he} orcs of the Violet Steppe do not separate magic from survival. A bone-singer who calls a storm is no different from a scout who reads the grass. Both are necessary. Both can get you killed.

\subsection*{What They Respect}
\begin{itemize}
    \item \textbf{Summoner (Ancestral):} The drum is the oldest Codex. A Ykrul summoner who can call a wolf-ancestor or a wind-horse is honored at any campfire.
    \item \textbf{Psion (Precognition):} The Ykrul call psions "Weather Readers" -- not because they predict rain, but because they see the shape of the future in the same way a rider sees the shape of the steppe.
    \item \textbf{Witch (Road craft):} A witch who can bind a promise with a knotted cord or curse a rival with a handful of grave-dust is a valued member of any clan.
\end{itemize}

\subsection*{What They Mistrust}
\begin{itemize}
    \item \textbf{Runekeeper:} Too slow. The steppe does not wait for you to carve a Rite.
    \item \textbf{Invoker:} Carrying Symbols from multiple Patrons is seen as divided loyalty. A Ykrul warrior serves one banner.
    \item \textbf{Cantor:} But only because a song that binds a crowd is a leash. The Ykrul value freedom too highly.
\end{itemize}

\subsection*{The Red Weather}
When a Ykrul witch curses a horse-thief, they do not use subtlety. They tie a red cord around a splinter of the thief's saddle and burn it in a storm. The "Red Weather" -- a sudden, localized squall -- is said to follow the cursed until they return what they stole. I have seen it happen twice. I have no explanation.

\section{Vhasia -- The Fractured Lens}

\lettrine{T}{he} Vhasians inherited the ruins of the Utaran Empire and learned to argue over the pieces. Their magic is legalistic, factional, and surprisingly bureaucratic.

\subsection*{What They Respect}
\begin{itemize}
    \item \textbf{Runekeeper:} A Vhasian Runekeeper is a notary with teeth. Their Codex contains Rites, yes, but also contracts, property deeds, and the terms of three unresolved feuds.
    \item \textbf{Invoker:} Multiple Patrons? In Vhasia, that is called "diversification." A noble who keeps Symbols from Aveh (for escape), Mykkiel (for contracts), and the Witness (for truth) is merely prudent.
    \item \textbf{Witch (Cunning craft):} Every Vhasian village has a hedge-witch who speaks for the old saints. The Church of the Everflame pretends not to notice.
\end{itemize}

\subsection*{What They Mistrust}
\begin{itemize}
    \item \textbf{Free Caster:} Too unpredictable. The Vhasian legal system cannot bill a free caster for damages if there is no contract.
    \item \textbf{Psion:} Reading minds without permission is a crime -- and the Vhasian Parlement has written three laws against it. Enforcement is spotty.
\end{itemize}

\subsection*{The Parlement of Dust}
In the catacombs beneath the Sun Palace, a secret court of Runekeepers adjudicates disputes between magical factions. They meet by candlelight and seal their judgments with wax mixed with grave-dust. No one admits the court exists. Everyone with real power obeys its rulings.

\section{Viterra -- The Hedge-Law Kingdom}

\lettrine{T}{he} Viterrans are practical people. Their magic grows in hedgerows and hides in parish ledgers. A witch in Viterra is not a mystic -- she is a "map adjuster" who knows where the boundary stones should really be.

\subsection*{What They Respect}
\begin{itemize}
    \item \textbf{Witch (Threshold):} The Viterran "hedge-wright" is a legal as well as magical profession. They draw property lines, witness oaths, and ensure that the new mill does not cross the ancient right-of-way.
    \item \textbf{Runekeeper:} But only those who specialize in "precedent Rites" -- spells that make a past judgment binding on the present.
    \item \textbf{Cantor:} A village Cantor who sings the old ballads at the harvest festival is beloved. One who sings politics is exiled.
\end{itemize}

\subsection*{What They Mistrust}
\begin{itemize}
    \item \textbf{Summoner:} Spirits cannot own property. The Viterran courts have ruled that a ghost's testimony is inadmissible.
    \item \textbf{Psion:} The Queen's Justiciars have a standing bounty on unlicensed telepaths. Privacy is the foundation of property law.
\end{itemize}

\subsection*{The Tilt-Yard Audits}
Before a Viterran Runekeeper can invoke a Rite that affects land ownership, they must "audit" the spell before a panel of Dawn-Knights. The knights strike the Codex with a sword. If the Rite survives three blows, it is lawful. I have seen a Codex shatter on the second strike. The Runekeeper was fined for "magical noise pollution."

\section{Ecktoria -- The Marble Ledger}

\lettrine{E}{cktoria} sits on the bones of the old Utaran Empire and pretends they are a throne. Their magic is ceremonial, political, and obsessed with appearances.

\subsection*{What They Respect}
\begin{itemize}
    \item \textbf{Runekeeper (Everflame):} The state religion has its own Runekeepers -- the "Censer-Knights" -- who march in processions and purify heretics with blessed fire.
    \item \textbf{Invoker:} Ecktorian nobles collect Symbols the way they collect statues of forgotten emperors. Display is more important than use.
    \item \textbf{Cantor:} A Cantor who can make a crowd weep with a hymn is a political asset. The Grand Magistrate employs three.
\end{itemize}

\subsection*{What They Mistrust}
\begin{itemize}
    \item \textbf{Witch:} Hedge-craft is "peasant magic." The Everflame preaches against it, but every noble's wife knows a witch who can brew a love potion.
    \item \textbf{Psion:} The Chain-Lanterns hunt psions with particular zeal. A mind that cannot be seen is a mind that might be planning rebellion.
    \item \textbf{Summoner:} The Ecktorian view on spirits is simple: if it is dead, it should stay dead. The Inquisition has burned three summoners in my lifetime.
\end{itemize}

\subsection*{The Censor's Mask}
At the annual Triumph, the Censor of Ecktoria wears a mask that is said to be a psionic artifact from the Great Silence. It lets him see the true intentions of every petitioner. I knelt before him once. He saw that I wanted his wine. He laughed and gave me a cup. The mask is real.

\section{Thepyrgos -- The City of Stairs}

\lettrine{T}{hepyrgos} is a nation of clerks and philosophers who believe that a properly argued case can bend reality. Their magic is academic, verbal, and exhausting.

\subsection*{What They Respect}
\begin{itemize}
    \item \textbf{Free Caster:} The University of Thepyrgos teaches the TAGS system as "practical metaphysics." A free caster with a degree can charge triple.
    \item \textbf{Psion:} The Synod has a standing committee on "mental hygiene." Licensed psions work as legal investigators. Unlicensed ones are hunted.
    \item \textbf{Witch (Bell craft):} The Thepyrgos bell-towers are thresholds. Bell-wardens are respected civic employees.
\end{itemize}

\subsection*{What They Mistrust}
\begin{itemize}
    \item \textbf{Runekeeper:} Too reliant on a single Patron. The Thepyrgos prefer negotiators, not worshippers.
    \item \textbf{Cantor:} A crowd that sings together is a crowd that might riot. The Archon limits public performances by Cantors to three per week.
\end{itemize}

\subsection*{The Silent Stair}
Behind the Synod Hall, a staircase that leads nowhere -- but psions claim it leads to the space between thoughts. The University paid me a handsome sum to climb it and report back. I climbed for an hour. The stairs did not end. I turned around and walked down. I was back in the same courtyard. I do not climb it anymore.

\section{Lethai -- The Fair Folk's Echo}

\lettrine{T}{he} elves of the Valewood do not practice magic. They -- and this is important -- \textit{are} magic, or at least they have forgotten where they end and the Weave begins. A Lethai Runekeeper is a contradiction. A Lethai Cantor is a redundancy. A Lethai Witch is a priestess of a religion that has no name.

\subsection*{What They Respect}
\begin{itemize}
    \item \textbf{Witch (Silk and Shed):} The Lethai have two witch traditions: the Silk Vigil (Inaea, mercy, continuity) and the Shed Vigil (Isoka, transformation, hard choices). Both are woven into the fabric of Lethai law.
    \item \textbf{Summoner:} Spirits are just neighbors who live in a different room. The Lethai treat them with courtesy -- and a healthy dose of suspicion.
    \item \textbf{Cantor:} A Lethai shira (song-weaver) can wrap a promise in a melody so beautiful that you forget you are bound. Their betrayal is legendary.
\end{itemize}

\subsection*{What They Mistrust}
\begin{itemize}
    \item \textbf{Invoker:} Why borrow power from a Patron when you could simply \textit{be} the power? The Lethai see Invokers as beggars in silk robes.
    \item \textbf{Psion:} Reading a Lethai mind without context is impossible -- their thoughts are saturated with memory, season, and relationship. Psions who try often leave mad.
    \item \textbf{Free Caster:} Too crude. The Lethai prefer their magic to have rules, even if the rules are secret.
\end{itemize}

\subsection*{The Ninth Walk}
Once a generation, a Lethai witch walks the "Ninth Walk" -- a pilgrimage through thresholds that do not exist on any map. She returns with a new law for the Lethai courts. The last Ninth Walk produced the rule: "No bell may ring nine times." No one remembers why. Everyone obeys.

\section{The Gray Wanderer's Conclusion}

I have been called a Runekeeper in Aeler, a witch in Aelaerem, a Cantor in Silkstrand, a Summoner in Zakov, a Psion in Thepyrgos, and a fool everywhere. The truth is simpler: every culture sees magic as it sees itself. The Aeler see a ledger. The Linn see an oath. The Thepyrgos see a debate. The Lethai see a dream.

You will find your own path. You will be misunderstood. You will be feared, honored, hunted, or hired -- sometimes all in the same day. When that happens, remember this: the power is not in the Rite or the Song or the silent thought. The power is in the choice to use them.

Now close the book. Go make some mistakes. I will be at the inn.

\vspace{0.5cm}
\hfill --- \textit{The Gray Wanderer}


\chapter{The Three Paths -- Lore and Roleplay}

\section{The Runekeeper -- Scholar of the Inked Vow}
\textit{"I do not beg. I record."}

Runekeepers are archivists of the divine. Their power comes from knowing the \textit{true name} of a Rite, the correct gesture, the exact syllable, the precise weight of the offering. They carry a \textbf{Codex} -- a personal grimoire -- and a \textbf{Thiasos} (familiar), a fragment of their Patron's attention made flesh.

\subsection{Lore -- Stone and Bark}
In Aeler, Runekeepers are \textit{Stone‑Scribes}, carving Rites into vault walls with their own blood. In Valewood, they weave Rites into living bark. Every Runekeeper knows that forgetting a single line can kill you.

\subsection{Roleplay Hooks}
\begin{itemize}
    \item Your Codex is incomplete -- the last page was torn out. What was on it?
    \item Your Thiasos is a sarcastic rat that only speaks in legal jargon.
    \item A rival Runekeeper claims you stole a Rite from their Codex. They want it back -- or your blood.
\end{itemize}

\subsection{Codex \& Thiasos in Play}
\begin{itemize}
    \item \textbf{Losing your Codex:} You cannot access any Rite until you recover it or spend 2 XP to recreate it from memory (requires a day of quiet and a Lore+Spirit test, DV 4).
    \item \textbf{Losing your Thiasos:} You lose the Patron's Gift (Imbuement) until you perform a ritual of re‑attunement.
\end{itemize}

\section{The Invoker -- Borrower of Stolen Fire}
\textit{"I don"t serve. I rent."}

Invokers carry \textbf{Symbols} -- physical anchors to Patrons they never fully swore to. They are gamblers, thieves, and desperate souls. In Silkstrand they are called \textit{Bridge‑Runners}. In Acasia, hedge‑witches with a dozen fetishes in their coat.

\subsection{Lore -- The Cracked Seal}
Every Invoker knows that a Symbol can break. When it does, the Patron might come to collect in person -- or send something smaller, sharper, and hungrier.

\subsection{Roleplay Hooks}
\begin{itemize}
    \item One of your Symbols was stolen by a former student. Now they use it against you.
    \item You have never met your Patron -- you found the Symbol in a ruin. Does the Patron even know you exist?
    \item You carry four Symbols (the limit before friction). Every morning, roll a die; on a 1, two Symbols clash. Mark 1 Fatigue.
\end{itemize}

\subsection{Symbols and Obligation}
A \textbf{well‑maintained Symbol} (polished, offered to weekly) reduces Obligation by 1 when you use it (once per session). A \textbf{compromised Symbol} (cracked, dull) can still be used with \textit{Crack the Seal}, but the GM may add a thematic curse.

\section{The Cantor -- Voice of the Unending Refrain}
\textit{"You think you need a lute? My larynx is older than any tree. Hum, and the world will listen."}

Cantors sing power into being. They are the most public of the paths -- you cannot hide a hymn. Their magic is as natural as breathing, yet as addictive as poppy wine.

\subsection{Lore -- The Body Remembers}
Cantors carry no Codex. No Thiasos. Their instrument is their own flesh. A Cantor who loses their voice cannot cast Rites -- but they can hum through clenched teeth, tap a rhythm on their ribs, or whistle. The body adapts.

\subsection{Roleplay Hooks}
\begin{itemize}
    \item You once sang a forbidden note, and now a fragment of a dead Cantor lives in your head.
    \item Your voice is unmistakable -- enemies recognise you by your timbre.
    \item You are trying to write a new Song, but every attempt ends in a headache. The Song is writing \textit{you}.
\end{itemize}

\section{The Summoner -- Pact-Whisperer}
\textit{"I do not command the dead. I ask. And sometimes -- when the price is right -- they answer."}

Summoners call spirits from the deep stone, the salt foam, the ashes of a pyre. Every culture has its own way -- contracts, offerings, drums, salt circles, cursed coins -- but the heart is the same: you need something the dead or the never‑born can give.

\subsection{Lore -- The Door and the Key}
In Aeler, summoners carve contracts into slate. In Ykrul camps, bone‑singers drum ancestors from the earth. In Zakov, fathom‑callers swallow cursed coins and whisper drowned names. The spirits remember. They have long memories.

\subsection{Roleplay Hooks}
\begin{itemize}
    \item You once called a spirit that still follows you -- not hostile, just \textit{curious}.
    \item Your mentor disappeared during a summoning. You found their Symbol but not their body.
    \item A spirit you bound years ago now demands release -- on its own terms.
\end{itemize}

\subsection{The Leash and the Bond}
When you summon, you establish a \textbf{Leash} -- the spiritual strain of keeping an Outsider in the world. Each time the spirit takes harm, acts against its nature, or crosses a ward, the Leash ticks. When it fills, the spirit acts once to its nature, then departs -- or turns hostile.

\clearpage

\chapter{Obligation -- The Debt You Carry}
\epigraph{"Every Rite is a receipt. Every Patron is a creditor. Pay on time, or they"ll repossess your soul."}{\textit{--- The Gray Wanderer}}

Obligation is \textbf{not a punishment}. It is a story engine. When your Obligation exceeds \texttt{Spirit + Presence}, you start marking Fatigue. When it doubles that cap, a \textbf{Patron Intrusion} occurs.

\section{Clearing Obligation -- General Ways}
\begin{itemize}
    \item \textbf{Act of Service (Downtime):} Perform a task aligning with the Patron's domain. Reduce Obligation by 1--2.
    \item \textbf{Ritual Purification (Scene):} Spend a scene in prayer or offering. Test \texttt{Lore + Spirit} (DV 3). Success reduces Obligation by 1.
    \item \textbf{Patron's Largess (GM Fiat):} When you dramatically fulfil the Patron's unstated goal, the GM may clear all Obligation.
\end{itemize}

\section{Patron Intrusions -- Narrative Consequences}
When Obligation exceeds \texttt{2 × (Spirit+Presence)}, the GM triggers an intrusion. Below are examples. The GM should tailor them to the current scene.

\clearpage

% ======================================================================
% PATRONS SECTION -- COMBINED FROM BOTH SOURCES
% ======================================================================

% Chapter 4: Patrons of the Gray Wanderer
% Expanded with categorization, cultural variants, and Gift options

\chapter{Patrons of the Gray Wanderer}

\epigraph{"You wish to walk the road of power? Then listen well. The world is old, and older still are the voices beneath it. We call them Patrons, though they were never sworn to us."}{\textit{--- The Gray Wanderer}}

Patrons are vast intelligences -- not gods, though some worship them as such. They are embodiments of concepts and forces. Each offers power, but always with cost: fatigue, scars upon fate, or a slow unweaving of one's own story.

This chapter presents twelve Patrons chosen for their utility and distinct themes, organized into four categories. Each entry includes the Patron's lore, cultural variants across the peoples of Fate's Edge, suggested Obligation acts, rites, and optional Patron's Gift variations beyond the standard Imbuement.

\section{The Wild -- Patrons of Nature and Season}

\subsection{Grimmir -- The Wild Speaker}
\label{patron:grimmir}

\textit{"In the space between planted row and forest edge, where the first grain met the first acorn, Grimmir waits. He speaks in the rustle of leaves, the migration of birds, and the patient growth of ancient oaks."}

Grimmir walks the threshold between cultivated lands and untamed wilderness. Neither wholly beast nor entirely human, he speaks in root and stone, seasons and sap.

\textbf{What They Represent:} Primal wisdom, seasonal cycles, the treaty between civilisation and wilderness.

\textbf{Cultural Variants:}
\begin{itemize}
    \item \textbf{Aeler (Dwarves):} Called \textit{the Root-Father}, patron of deep fungus and stone-lichen. His shrines are carved into living rock where moss is never cleared.
    \item \textbf{Ykrul (Orcs):} A sky‑wolf who herds the seasons; shamans invoke Grimmir to predict the spring thaw. His symbol is a wolf's paw print in ash.
    \item \textbf{Aelaerem (Halflings):} \textit{the Thorn-Father}, keeper of hedgerow magic. Villagers tie red thread around his roadside stones before planting.
    \item \textbf{Vilikari (Federates):} A two‑faced figure -- one side a sower, one side a reaper -- invoked at both planting and harvest festivals.
\end{itemize}

\textbf{Playing a Follower of Grimmir:}
\textbf{Strategy:} Grimmir's rites reward patience, observation, and environmental manipulation. You are a controller -- slowing enemies with thorns, creating advantageous terrain.

\textbf{Suggested Acts of Obligation:}
\begin{itemize}
    \item Plant a tree or restore a damaged grove.
    \item Kill a predator that has grown unnaturally bold.
    \item Teach a villager to read a seasonal sign.
\end{itemize}

\textbf{Patron's Gift (Imbuement):} Your Thiasos or your own hands glow with verdant energy. For one scene, you gain +1 die to Survival and Nature rolls, and your melee strikes deal +1 Harm against creatures of corruption or blight. This replaces the standard +1 Melee / +1 Thematic -- you choose which benefit applies when you use the Gift.

\textbf{Alternative Gift -- Wildshape (Minor Transformation):} You may spend an action and mark 1 Obligation to assume the form of a small beast (wolf, hawk, badger, or similar) for up to one scene. While transformed, you gain +2 dice to Stealth and Perception in natural terrain, but cannot speak or use tools. You may revert as a free action. This counts as your Patron's Gift for the scene -- you cannot also gain the Imbuement bonuses.

\textbf{Rites -- Strategic Use:}
\textbf{Low Rites}
\begin{itemize}
    \item \textit{The Speaking Seed} (4 XP): Commune with local plant life.
    \item \textit{The Season's Turn} (5 XP): Align with seasonal power.
\end{itemize}
\textbf{Standard Rites}
\begin{itemize}
    \item \textit{Verdant Tongue} (8 XP): Speak with all plant and animal life in Near range.
    \item \textit{The Thornveil} (7 XP): Raise a living barrier.
\end{itemize}
\textbf{Advanced Rites}
\begin{patronbox}{The World's Wound}
\textbf{Cost:} 13 XP, Obligation 7 segments \\
\textbf{Effect:} Heal a blighted place or quicken a damaged ecosystem. Clear taint, restore fertility, awaken dormant forces. Begin an \textit{Ecosystem Restoration} [8] clock.
\end{patronbox}

\textbf{GM Guidance -- Grimmir's Intrusions:}
\begin{enumerate}
    \item A tree you passed yesterday is now blocking your path -- with a face carved into it that looks like yours.
    \item All food you touch spoils within an hour unless you offer the first bite to the earth.
    \item A pack of wolves shadows the party for a full day. They do not attack -- they simply watch.
\end{enumerate}

\clearpage

\subsection{Raéyn -- Mistress of the Sea}
\label{patron:raeyn}

\textit{"Mark the tide, name your course, and trust the wave‑road. But speak ill of Khemesh, and even I may let the deep take you."}

Raéyn is the tempestuous goddess of the sea, the restless tide that carries news between shores. Mother to all who sail, her voice the wind that fills sails and her moods the storms that test every mariner's resolve.

\textbf{What They Represent:} Tides, change, sea travel, storms, the dual nature of the ocean.

\textbf{Cultural Variants:}
\begin{itemize}
    \item \textbf{Kahfagia:} \textit{the Lantern‑Tide}. Pilots light a beacon to Raéyn before every night voyage. Her symbol is a wave curling into a spiral.
    \item \textbf{Linn:} \textit{the Whale‑Mother}. Raéyn is the great she‑whale who carries the drowned to the afterlife in her belly. Funerals at sea involve a song and a single white stone.
    \item \textbf{Zakov:} \textit{the Salt‑Queen}. Pirates carve her face on their figureheads -- one eye open (fair seas), one eye closed (storms ahead).
\end{itemize}

\textbf{Playing a Follower of Raéyn:}
\textbf{Strategy:} You are the party's navigator and weather‑reader. Raéyn's rites excel at travel on or near water, weather manipulation, and the ebb and flow of tides.

\textbf{Suggested Acts of Obligation:}
\begin{itemize}
    \item Pour a libation into the sea before a voyage, asking for nothing in return.
    \item Rescue a drowning creature -- a sailor, a fisher, even a seabird.
    \item Walk the shore at dawn or dusk and listen. If you hear the waves speaking a name, speak it back.
\end{itemize}

\textbf{Patron's Gift (Imbuement):} Your touch carries the scent of brine and rolling waves. For one scene, you gain +1 die to Athletics (swimming) and Navigation, and you ignore one level of environmental penalty from storms or choppy water.

\textbf{Alternative Gift -- Tide‑Step:} Once per session as a free action, you may treat any water‑based difficult terrain as stable ground for one exchange. Alternatively, you may move across the surface of water as if it were solid for up to Near distance. This counts as your Patron's Gift activation.

\textbf{Rites -- Strategic Use:}
\textbf{Low Rites}
\begin{itemize}
    \item \textit{Tidemark's Blessing} (4 XP): Treat slick or water-slicked footing as stable.
    \item \textit{Whispering Currents} (5 XP): Learn the safest near‑term route across water.
\end{itemize}
\textbf{Standard Rites}
\begin{itemize}
    \item \textit{The Changing Tide} (8 XP): Bias currents and water levels in a zone.
    \item \textit{Wave‑Road Blessing} (7 XP): Consecrate a wave‑road for +2 dice to travel.
\end{itemize}
\textbf{Advanced Rites}
\begin{patronbox}{The Storm‑Queen's Hand}
\textbf{Cost:} 13 XP, Obligation 7 segments \\
\textbf{Effect:} Shape a storm‑band over a zone (sea or shore). Choose two modes: propulsion, concealment, or smite.
\end{patronbox}

\textbf{GM Guidance -- Raéyn's Intrusions:}
\begin{enumerate}
    \item A sudden fog rolls in -- all navigation rolls at --2 dice for the scene.
    \item The tide turns early, stranding a vessel you need.
    \item A dolphin, seal, or sea eagle follows the party for a day. It will deliver one short message to the nearest follower of Raéyn if fed something shiny.
\end{enumerate}

\clearpage

\subsection{The Carrion-King -- Master of Endings}
\label{patron:carrionking}

\textit{"What crumbles feeds what grows. What dies becomes the soil of tomorrow's triumph. Do not mourn the fallen -- harvest them."}

The Carrion‑King is not a lord of rot for its own sake. He is the alchemist of endings, the patient gardener who sees every corpse as compost.

\textbf{What They Represent:} Decay, renewal, transformation.

\textbf{Cultural Variants:}
\begin{itemize}
    \item \textbf{Vhasia:} \textit{the Dust‑Accountant}. Soldiers leave a button from a fallen comrade on a cairn -- the King will turn that button into a bullet for the next battle.
    \item \textbf{Aeler:} \textit{the Under‑Gardener}. Dwarven crypts have bioluminescent fungus fed by the dead. Tending it is an act of devotion.
    \item \textbf{Inquisitor Prime's View:} Witch‑hunters see the Carrion‑King as a necromantic menace, confusing "natural decay" with "unholy resurrection".
\end{itemize}

\textbf{Playing a Follower of the Carrion‑King:}
\textbf{Strategy:} You are a harvester of potential. In combat, you weaken enemies with rot and decay, then use their defeat to empower your allies.

\textbf{Suggested Acts of Obligation:}
\begin{itemize}
    \item Compost a fallen enemy's body and use the soil to plant something edible.
    \item Salvage a ruined structure and give the reclaimed materials to a needy family.
    \item Teach someone that loss is not failure.
\end{itemize}

\textbf{Patron's Gift (Imbuement):} Your touch carries the scent of turned earth and gentle decay. For one scene, you gain +1 die to Medicine and Craft rolls involving organic materials, and you may convert one minor item of personal value into a single dose of natural remedy or poison.

\textbf{Alternative Gift -- Harvest Resonance:} Once per scene after an enemy is defeated or a significant ending occurs, you may draw a sliver of that ending's energy. Gain +1 die on your next action, or clear 1 Fatigue. This counts as your Patron's Gift for the scene.

\textbf{Rites -- Strategic Use:}
\textbf{Low Rites}
\begin{itemize}
    \item \textit{Rite of Consuming Rot} (5 XP): Accelerate natural decay to weaken organic materials.
    \item \textit{Rite of the Harvested End} (4 XP): Extract value from a recent ending.
\end{itemize}
\textbf{Standard Rites}
\begin{itemize}
    \item \textit{Rite of the Fertile Death} (8 XP): Transform death into growth -- create advantageous terrain or grant healing bonuses.
    \item \textit{Rite of the Transformed Spirit} (7 XP): Channel the essence of a recently deceased being.
\end{itemize}
\textbf{Advanced Rites}
\begin{patronbox}{The Eternal Cycle}
\textbf{Cost:} 14 XP, Obligation 8 segments \\
\textbf{Effect:} Complete a transformation cycle. Destroy one major asset, enemy, or obstacle, and create something new of equal or greater value.
\end{patronbox}

\textbf{GM Guidance -- Carrion‑King's Intrusions:}
\begin{enumerate}
    \item A dead creature you passed earlier is now following you -- not as an enemy, but as a witness.
    \item All food in the party's packs spoils instantly, except for a single loaf of black bread.
    \item A villager you saved now bears a small rot‑mark that spreads one inch each day.
\end{enumerate}

\clearpage

\section{The Shadow -- Patrons of Secrets and Transformation}

\subsection{Ikasha -- She Who Sleeps}
\label{patron:ikasha}

\textit{"Blow out the candle. If the room listens back, ask softly. At the next crossroads, the raven waits -- and the shadow remembers your passing."}

Ikasha is the hush between footfalls, the patience of dark water, the black‑feathered watcher at every threshold. In stillness she gathers what might be.

\textbf{What They Represent:} Latent potential, shadow, thresholds, forgotten things.

\textbf{Cultural Variants:}
\begin{itemize}
    \item \textbf{Mistlands:} \textit{the Bell‑Shrouded}; bell‑wardens leave a single rope unknotted for her, so she may ring silence when needed.
    \item \textbf{Zakov:} \textit{the Tide Watcher}, a smuggler's patron. She rises with the night tide, and her ravens carry messages between pirate captains.
    \item \textbf{Aeler:} \textit{the Deep Listener}; miners whisper to her before descending into new shafts.
\end{itemize}

\textbf{Playing a Follower of Ikasha:}
\textbf{Strategy:} Ikasha's rites reward stealth, deception, and strategic patience. You are the one who waits for the guard's shift change.

\textbf{Suggested Acts of Obligation:}
\begin{itemize}
    \item Spend a night in absolute darkness, speaking no word and lighting no flame.
    \item Carry a secret for someone -- a truth that would ruin them -- and tell no one for a full season.
    \item At a crossroads, leave a black feather and a single coin. Do not look back.
\end{itemize}

\textbf{Patron's Gift (Imbuement):} Your form seems to drink the light. For one scene, you gain +1 die to Stealth and Deception, and your footsteps make no sound unless you choose otherwise.

\textbf{Alternative Gift -- Shadow Meld:} Once per session as an action, you may step into any shadow larger than yourself and emerge from another shadow within Near range. You cannot be targeted while moving between shadows. This counts as your Patron's Gift activation.

\textbf{Rites -- Strategic Use:}
\textbf{Low Rites}
\begin{itemize}
    \item \textit{Touch the Umbral Veil} (4 XP): Start Controlled on one Stealth roll.
    \item \textit{Rite of the Crossroads Raven} (5 XP): Summon an omen‑raven for +1 die to navigation or pursuit.
\end{itemize}
\textbf{Standard Rites}
\begin{itemize}
    \item \textit{Draw from the Umbral Reservoir} (8 XP): Grant +2 dice to stealth, deception, or resolve.
    \item \textit{Secret Keeper's Burden} (9 XP): Compel a truthful answer to one direct question.
\end{itemize}
\textbf{Advanced Rites}
\begin{patronbox}{Become the Shadow at the Crossroads}
\textbf{Cost:} 12 XP, Obligation 7 segments \\
\textbf{Effect:} Intangible to mundane harm, pass through thresholds, +2 dice to Stealth, auto‑succeed one escape.
\end{patronbox}

\textbf{GM Guidance -- Ikasha's Intrusions:}
\begin{enumerate}
    \item A raven lands on your shoulder and whispers a secret you never wanted to know.
    \item Your shadow stretches into the shape of a woman with a crown of black feathers.
    \item A door you passed through earlier is now locked from the other side -- and the key is in your pocket.
\end{enumerate}

\clearpage

\subsection{Isoka -- Angel of Serpents}
\label{patron:isoka}

\textit{"Do not mourn the skin you shed. It was never meant to last. The venom that burns away the old self is the same that grants the strength to become new."}

Isoka teaches that every self is temporary -- identity is a skin to be shed so growth can continue. Her serpents are omens and teachers.

\textbf{What They Represent:} Transformation, renewal, venom, the casting off of old selves.

\textbf{Cultural Variants:}
\begin{itemize}
    \item \textbf{Zakov:} \textit{the Tide‑Shedder}. Pirates leave a lock of hair and a written name at her sunken temple; eels eat both.
    \item \textbf{Acasia:} \textit{the Copperhead}. Hedge‑cultists shed identities to infiltrate noble houses.
    \item \textbf{Inquisitor Prime's View:} Isoka's followers are shape‑shifters and liars, abominations who erase themselves to avoid judgment.
\end{itemize}

\textbf{Playing a Follower of Isoka:}
\textbf{Strategy:} You are the party's infiltrator and agent of painful renewal. Isoka's rites excel at disguise, poison, escape, and breaking bonds.

\textbf{Suggested Acts of Obligation:}
\begin{itemize}
    \item Shed a physical token of your old life -- a piece of clothing, a weapon, a letter -- and leave it at a crossroads.
    \item Forgive someone who wronged you, publicly and completely, as an act of transformation.
    \item Help someone else change -- cut their hair, give them new clothes, teach them a new name.
\end{itemize}

\textbf{Patron's Gift (Imbuement):} Your movements become fluid and sinuous. For one scene, you gain +1 die to Athletics (climbing, squeezing) and Deception, and you may re-roll one failed escape attempt.

\textbf{Alternative Gift -- Shed Skin:} Once per session when you are Bound, Grappled, or suffering from a Condition that restricts movement, you may spend an action and mark 1 Obligation to shed that restriction. Describe how you slip free like a snake shedding its skin. This counts as your Patron's Gift activation.

\textbf{Rites -- Strategic Use:}
\textbf{Low Rites}
\begin{itemize}
    \item \textit{Venomous Benediction} (5 XP): Bless an ally's strike with +1 Harm and possible Poison.
    \item \textit{Loosening Skin} (4 XP): +1 die to resist an ongoing Condition.
\end{itemize}
\textbf{Standard Rites}
\begin{itemize}
    \item \textit{The Shedding} (8 XP): +2 dice to resist a named Condition; declare a minor contingency retroactively.
    \item \textit{Forked Tongue} (7 XP): Ambiguous persuasion -- success may generate leverage instead of SB.
\end{itemize}
\textbf{Advanced Rites}
\begin{patronbox}{Complete Metamorphosis}
\textbf{Cost:} 13 XP, Obligation 7 segments \\
\textbf{Effect:} Change appearance, voice, and mannerisms completely. Warded against recognition.
\end{patronbox}

\textbf{GM Guidance -- Isoka's Intrusions:}
\begin{enumerate}
    \item A piece of your old self returns -- a letter from a dead lover, a scar you had forgotten.
    \item Your shadow moves independently, pointing at people who are about to change -- to leave, to betray, to die.
    \item You begin to shed skin overnight. If you burn the cast, you forget one minor memory. If you keep it, someone will steal it.
\end{enumerate}

\clearpage

\subsection{Aveh -- The Rider Behind the Storm}
\label{patron:aveh}

\textit{"I am the track the storm devours. I am the freedom you fear and the silence you crave. To ride with me is to be unbound -- and to be unbound is to vanish."}

Aveh is the faceless rider at the horizon's edge, offering liberty at the cost of belonging. Wanderers, exiles, oath‑breakers, and rebels whisper Aveh's name.

\textbf{What They Represent:} Freedom, erasure, storms, exile.

\textbf{Cultural Variants:}
\begin{itemize}
    \item \textbf{Ykrul:} A sky‑spirit that rides with winter gales; shamans invoke Aveh to hide a dying person's tracks from the dead.
    \item \textbf{Linn:} A sea‑fog rider; sailors whisper Aveh's name when they must slip through a blockade unseen.
    \item \textbf{Acasia:} The "dust‑walker" who appears at crossroads; followers leave a coin on the lintel when they flee a blood feud.
\end{itemize}

\textbf{Playing a Follower of Aveh:}
\textbf{Strategy:} Aveh's rites excel at escape, stealth, and erasure. You are not a front‑line fighter -- you are the one who slips away.

\textbf{Suggested Acts of Obligation:}
\begin{itemize}
    \item Cut a binding contract and scatter the pieces to the wind.
    \item Help someone flee oppression -- a slave, a debtor, a threatened child.
    \item Spend a night sleeping under open sky with no shelter.
\end{itemize}

\textbf{Patron's Gift (Imbuement):} Your presence becomes hard to hold in memory. For one scene, you gain +1 die to Stealth and Escape, and anyone who sees you must succeed on a Resolve test (DV 2) to remember your face clearly.

\textbf{Alternative Gift -- Storm‑Step:} Once per session as a move action, you may vanish in a swirl of wind and dust and reappear anywhere within Near range that you can see. This movement does not provoke opportunity attacks. This counts as your Patron's Gift activation.

\textbf{Rites -- Strategic Use:}
\textbf{Low Rites}
\begin{itemize}
    \item \textit{Storm-Step} (4 XP): Slip free of notice or restraint for a scene.
    \item \textit{Exile's Banner} (5 XP): Mark a target as "outside" -- socially erased.
\end{itemize}
\textbf{Standard Rites}
\begin{itemize}
    \item \textit{Erase the Road} (8 XP): Trails vanish, records blur, pursuers lose the way.
    \item \textit{The Rider's Mark} (9 XP): Bestow Aveh's mark on an ally -- they resist capture but become forgettable.
\end{itemize}
\textbf{Advanced Rites}
\begin{patronbox}{Storm's Refuge}
\textbf{Cost:} 12 XP, Obligation 6 segments \\
\textbf{Effect:} Shroud your company -- unseen or forgotten by pursuers. All present mark 1 Corruption or forget one bond at the end.
\end{patronbox}

\textbf{GM Guidance -- Aveh's Intrusions:}
\begin{enumerate}
    \item A storm erases your tracks -- and also the memory of your face from the last town.
    \item Your shadow detaches and wanders off for an hour.
    \item A forgotten debt suddenly resurfaces -- someone you wronged now knows exactly where you are.
\end{enumerate}

\clearpage

\section{The Covenant -- Patrons of Law and Truth}

\subsection{The Inquisitor Prime -- Purity and Domination}
\label{patron:inquisitor}

\textit{"Mercy is weakness. Purity is survival. Where corruption hides, we carve it out."}

Among zealots and witch‑hunters, the Inquisitor Prime is venerated as the hand of absolute purity and the sword of uncompromising order.

\textbf{What They Represent:} Purity, order, judgment, eradication of the tainted.

\textbf{Cultural Variants:}
\begin{itemize}
    \item \textbf{Ecktoria (Everflame):} \textit{the Unquenchable Fire}. Priests carry censers of perpetual coal; rites focus on purification and judgment.
    \item \textbf{Aeler:} \textit{the Light of Forges}. Dwarven smiths kindle hearths with a brand from the temple's eternal flame.
    \item \textbf{The Inquisitor Prime's own cult:} Some sects believe the Prime is not a Patron but a mortal who ascended by burning away his own humanity.
\end{itemize}

\textbf{Playing a Follower of the Inquisitor Prime:}
\textbf{Strategy:} You are a hunter of the arcane and the aberrant. Your rites excel at revealing lies, dispelling illusions, and binding supernatural foes.

\textbf{Suggested Acts of Obligation:}
\begin{itemize}
    \item Burn a heretical text or destroy a cursed object.
    \item Expose a charlatan who uses magic to deceive the innocent.
    \item Hunt down a rogue summoner or corrupt Runekeeper.
\end{itemize}

\textbf{Patron's Gift (Imbuement):} Your gaze burns with righteous fire. For one scene, you gain +1 die to Insight and Intimidation when confronting supernatural creatures, and your melee attacks ignore 1 point of Armor from undead or demons.

\textbf{Alternative Gift -- Truth's Gaze:} Once per session as an action, you may fix your gaze on a target within Near range. For the next minute, they cannot tell a direct lie to you; they may evade or refuse to answer, but any statement they make must be factually true. This counts as your Patron's Gift activation.

\textbf{Rites -- Strategic Use:}
\textbf{Low Rites}
\begin{itemize}
    \item \textit{Rite of the Pure Flame} (4 XP): +2 dice to resist supernatural influence.
    \item \textit{Rite of the Unclouded Eye} (5 XP): +2 dice to spot illusions or magical deception.
\end{itemize}
\textbf{Standard Rites}
\begin{itemize}
    \item \textit{Rite of the Cleansing Light} (8 XP): Dispel a magical effect or weaken it.
    \item \textit{Rite of the Marked Prey} (7 XP): Mark a supernatural target -- they cannot hide with illusion.
\end{itemize}
\textbf{Advanced Rites}
\begin{patronbox}{The Final Admonition}
\textbf{Cost:} 14 XP, Obligation 8 segments \\
\textbf{Effect:} Attempt to annihilate one supernatural foe. Target tests Spirit+Resolve (DV 5). On failure, destroyed outright. On success, suffers Harm 3 and --2 dice this scene.
\end{patronbox}

\textbf{GM Guidance -- Inquisitor Prime's Intrusions:}
\begin{enumerate}
    \item An old investigation resurfaces -- someone you once cleared is now accused again.
    \item Your zeal attracts attention from both enemies and \textit{allies} who fear your methods.
    \item A supernatural entity you wounded escapes and seeks revenge.
\end{enumerate}

\clearpage

\subsection{The Witness -- Keeper of Uncomfortable Truths}
\label{patron:witness}

\textit{"I will show you what you would rather forget. But first, you must forget what you think you know."}

The Witness remembers what others bury. Every shadow cast and oath broken is a line in her unending ledger. She does not judge; she merely \textit{records}.

\textbf{What They Represent:} Truth, revelation, memory, the weight of knowledge.

\textbf{Cultural Variants:}
\begin{itemize}
    \item \textbf{Ecktoria:} \textit{the Ink‑Eye}. Censor's clerks whisper that the Witness copies every redacted line into a private scroll.
    \item \textbf{Vhasia:} \textit{the Fractured Mirror}. The Sun‑court's spies pray to her to reveal the true loyalty of rivals.
    \item \textbf{Inquisitor Prime's View:} The Witness and the Inquisitor are often confused -- but the Inquisitor destroys untruth; the Witness merely exposes it.
\end{itemize}

\textbf{Playing a Follower of the Witness:}
\textbf{Strategy:} You are the party's investigator and interrogator. Your rites excel at detecting lies and uncovering hidden information.

\textbf{Suggested Acts of Obligation:}
\begin{itemize}
    \item Expose a lie told by someone in power.
    \item Write down a secret you have never told anyone, then burn the paper.
    \item Testify truthfully in a legal proceeding, even if it hurts your cause.
\end{itemize}

\textbf{Patron's Gift (Imbuement):} Your eyes glimmer with remembered light. For one scene, you gain +1 die to Investigation and Insight, and you may ask one factual question about a scene or object; the GM must answer truthfully, though perhaps cryptically.

\textbf{Alternative Gift -- Recorded Echo:} Once per session, you may touch a physical trace (a footprint, a torn cloth, a drop of blood) and instantly know the last significant action taken by the person who left it (e.g., "they were running," "they hid something under the floor"). This counts as your Patron's Gift activation.

\textbf{Rites -- Strategic Use:}
\textbf{Low Rites}
\begin{itemize}
    \item \textit{Lingering Glimpse} (4 XP): +1 die to investigate something related to a trace you hold.
    \item \textit{Piercing Scrutiny} (5 XP): +1 die to detect deception in a consecrated zone.
\end{itemize}
\textbf{Standard Rites}
\begin{itemize}
    \item \textit{Echoing Truth} (8 XP): Compel a truthful answer to one direct question.
    \item \textit{Immutable Record} (7 XP): Bind an agreement. Breakers suffer reputation damage.
\end{itemize}
\textbf{Advanced Rites}
\begin{patronbox}{The Unveiled Heart}
\textbf{Cost:} 12 XP, Obligation 6 segments \\
\textbf{Effect:} Target suffers --2 dice to conceal true emotions or lies. Any successful social roll they make generates 1 SB.
\end{patronbox}

\textbf{GM Guidance -- The Witness's Intrusions:}
\begin{enumerate}
    \item You remember something you never witnessed -- a crime, a conversation, a death -- in perfect detail.
    \item A mirror shows not your face, but the face of someone you have wronged.
    \item A secret you have kept is now known to a stranger -- a blackmailer, a journalist, a rival.
\end{enumerate}

\clearpage

\subsection{The Traveler -- Lord of the Open Road}
\label{patron:traveler}

\textit{"One foot in a promise, and the road will meet you halfway. But break your word to the way, and the way will break you."}

The Traveler is not merely one who shows the path -- they \textit{are} the path, existing in the pause between steps and in the choice of which fork to take when roads diverge.

\textbf{What They Represent:} Ways, journeys, thresholds, safe passage.

\textbf{Cultural Variants:}
\begin{itemize}
    \item \textbf{Kahfagia:} \textit{the Lantern‑Bearer}. Pilots light a beacon at the start of every convoy.
    \item \textbf{Ykrul:} \textit{the Hoof‑Print}. Riders leave a single arrowhead at a crossroads for luck.
    \item \textbf{Acasia:} \textit{the Broken Wheel}. Desperate travellers scratch the Traveler's sigil -- a spiral -- on milestones.
\end{itemize}

\textbf{Playing a Follower of the Traveler:}
\textbf{Strategy:} You are the party's guide and logistician. Your rites excel at navigation, escape, and protecting caravans.

\textbf{Suggested Acts of Obligation:}
\begin{itemize}
    \item Walk a mile of road barefoot, carrying a stone in each hand.
    \item Offer shelter to a lost traveller for a night -- no questions, no payment.
    \item Repair a broken signpost, milestone, or footbridge.
\end{itemize}

\textbf{Patron's Gift (Imbuement):} Your feet never seem to tire. For one scene, you gain +1 die to Athletics (running, climbing) and Navigation, and you ignore one level of difficult terrain.

\textbf{Alternative Gift -- Shortcut:} Once per session when moving from one location to another (including during travel montages), you may declare a hidden shortcut that reduces the remaining travel clock by 2 segments. The GM may describe a minor cost or complication (e.g., "the path is steep -- mark 1 Fatigue"). This counts as your Patron's Gift activation.

\textbf{Rites -- Strategic Use:}
\textbf{Low Rites}
\begin{itemize}
    \item \textit{Road‑Sense} (4 XP): Unerringly pick the fastest safe route.
    \item \textit{Traveler's Boon} (5 XP): Ignore one level of difficult terrain.
\end{itemize}
\textbf{Standard Rites}
\begin{itemize}
    \item \textit{The Waymark} (8 XP): Declare a lane as permitted passage for your party.
    \item \textit{The Bridge Between} (7 XP): Relocate a willing target within Far range.
\end{itemize}
\textbf{Advanced Rites}
\begin{patronbox}{The Wayfarer's Covenant}
\textbf{Cost:} 14 XP, Obligation 8 segments \\
\textbf{Effect:} You and up to 5 travellers swear a covenant of safe passage. +2 Effect on travel rolls, re‑roll one failed navigation test per scene.
\end{patronbox}

\textbf{GM Guidance -- The Traveler's Intrusions:}
\begin{enumerate}
    \item A path you have walked before now leads somewhere else -- a different building, a different forest.
    \item Your boots wear out instantly. You must walk barefoot for a scene.
    \item A stranger approaches on the road and asks for directions. If you lie, you double your current Obligation.
\end{enumerate}

\clearpage

\section{The Infinite -- Patrons of the Unseen and Unknowable}

\subsection{Khemesh -- The Abyssal Maw}
\label{patron:khemesh}

\textit{"In the trench without light, the Maw waits. Even silence drowns."}

Khemesh is not merely a lord of the depths -- he is the hunger beneath them, a pressure older than seas. Those who bargain with him are marked by the abyss.

\textbf{What They Represent:} Depth, pressure, inevitability, eldritch terror.

\textbf{Cultural Variants:}
\begin{itemize}
    \item \textbf{Zakov:} \textit{the Drowned Prince}. Pirates toss a silver coin overboard before a raid -- if it sinks, Khemesh approves.
    \item \textbf{Linn:} \textit{the Reef‑Father}. Fishermen leave the first catch of the day to him, lest the tide turn and drag their nets into the abyss.
    \item \textbf{Mistlands:} \textit{the Bell‑Breaker}. When a bell line fails, wardens whisper that Khemesh has breathed on the clappers.
\end{itemize}

\textbf{Playing a Follower of Khemesh:}
\textbf{Strategy:} Khemesh's rites excel at control, crushing force, and psychological terror. You are a hammer -- slow, inevitable, and devastating.

\textbf{Suggested Acts of Obligation:}
\begin{itemize}
    \item Throw a valuable object into the deepest water you can reach.
    \item Spend a night submerged up to your neck in cold water.
    \item Reveal a secret you have kept for more than a year where water can hear.
\end{itemize}

\textbf{Patron's Gift (Imbuement):} Your presence feels like the deep ocean -- heavy, unyielding. For one scene, you gain +1 die to Intimidation and Endurance, and hostile creatures must test Resolve to move closer to you.

\textbf{Alternative Gift -- Crushing Silence:} Once per session as a reaction when a creature within Near range speaks or casts a spell, you may impose a bubble of oppressive silence around them for one exchange. They cannot speak, and any spell with a verbal component fails. This counts as your Patron's Gift activation.

\textbf{Rites -- Strategic Use:}
\textbf{Low Rites}
\begin{itemize}
    \item \textit{Whisper of the Trench} (4 XP): Target suffers --1 die on next action.
    \item \textit{Rite of Crushing Silence} (5 XP): Establish oppressive silence in a zone.
\end{itemize}
\textbf{Standard Rites}
\begin{itemize}
    \item \textit{Pressure of the Maw} (7 XP): Pin a target with invisible force -- treat as [ENTANGLE].
    \item \textit{Rite of the Abyssal Vision} (9 XP): See the world as Khemesh does -- fractured, alien.
\end{itemize}
\textbf{Advanced Rites}
\begin{patronbox}{The Maw Opens}
\textbf{Cost:} 12 XP, Obligation 8 segments \\
\textbf{Effect:} Reality folds inward. Enemies act at Desperate Position by default. Structures fracture under immense weight.
\end{patronbox}

\textbf{GM Guidance -- Khemesh's Intrusions:}
\begin{enumerate}
    \item Water pressure quadruples in the room -- wooden objects groan, glassware cracks.
    \item A nightmare tentacle manifests from a nearby body of water and steals a shiny object.
    \item The party begins to \textit{sink} -- not physically, but emotionally. Each character marks 1 Fatigue or confesses a secret fear aloud.
\end{enumerate}

\clearpage

\subsection{The Ninth -- Beyond Comprehension}
\label{patron:ninth}

\textit{"Eight figures bind the world; the Ninth overflows it. Drink, and remember that the cup forgets its shape before the song is done."}

Beyond the eighth figure of Sacred Geometry lies a rim no compass will hold. The Ninth is not absence but surplus: knowledge poured past the lip of any vessel, truth that overgrows the frame that seeks to bind it.

\textbf{What They Represent:} Hyperdimensional knowledge, incomprehensible truths, information overflow.

\textbf{Cultural Variants:}
\begin{itemize}
    \item \textbf{Thepyrgos:} \textit{the Lost Syllable}. Archivists believe the Ninth is a word erased from every dictionary -- but still whispered in dreams.
    \item \textbf{Aeler:} \textit{the Deep Count}. Geometers speak of a number that cannot be written, only approximated.
    \item \textbf{Inquisitor Prime's View:} The Ninth is the ultimate heresy -- knowledge that corrupts by its very existence.
\end{itemize}

\textbf{Playing a Follower of the Ninth:}
\textbf{Strategy:} You seek patterns others cannot see. Your rites grant impossible insights, but each revelation strains the mind.

\textbf{Suggested Acts of Obligation:}
\begin{itemize}
    \item Solve a puzzle or riddle that has stumped others for years.
    \item Destroy a book that contains a comfortable lie, replacing it with a painful truth.
    \item Teach someone a concept they are not ready to understand -- then watch what happens.
\end{itemize}

\textbf{Patron's Gift (Imbuement):} Your thoughts race ahead of your tongue. For one scene, you gain +1 die to Lore and Investigation, and you may ask one question about a pattern or hidden connection; the GM must answer with a cryptic truth.

\textbf{Alternative Gift -- Glimpse Beyond:} Once per session, you may look at a complex situation (a crime scene, a puzzle lock, a battle formation) and instantly intuit the single most important hidden detail. The GM gives you one piece of information that would otherwise require a successful Investigation roll. This counts as your Patron's Gift activation.

\textbf{Rites -- Strategic Use:}
\textbf{Low Rites}
\begin{itemize}
    \item \textit{Ninth Figure} (4 XP): +1 die to Lore or Investigation when wrestling complex systems.
    \item \textit{Flooded Scriptorium} (5 XP): Ask three questions upon any topic that leaves a written wake.
\end{itemize}
\textbf{Standard Rites}
\begin{itemize}
    \item \textit{Hypercognitive Lattice} (8 XP): Treat Wits as +2 higher for information-processing rolls.
    \item \textit{Breaking Glyph} (7 XP): Perfect internal grasp of one complex system for the session.
\end{itemize}
\textbf{Advanced Rites}
\begin{patronbox}{The Ninth Revelation}
\textbf{Cost:} 14 XP, Obligation 8 segments \\
\textbf{Effect:} Glimpse a shard of Ninth truth -- perfect command of a knotty system or comprehension of an "impossible" relation.
\end{patronbox}

\textbf{GM Guidance -- The Ninth's Intrusions:}
\begin{enumerate}
    \item You see patterns in everything -- the way dust falls, the arrangement of stones. It is useful, but exhausting. Mark 1 Fatigue.
    \item A fact you learned from the Ninth is true, but \textit{should not be}. Others find it disturbing when you speak it aloud.
    \item You begin to forget mundane things (names, faces, what you ate) while remembering impossible geometries with perfect clarity.
\end{enumerate}

\section{The Necessary Evils -- Patrons of Bond, Law, and Ruin}

Between the wild places and the shadowed thresholds, between the covenant of truth and the infinite abyss, there exist Patrons who answer not to hope or fear, but to \textit{necessity}. They are the ones mortals call when a vow must be witnessed, when a curse must be borne, or when a web must be woven. They are not kind. They are not cruel. They are \textit{inevitable}.

\subsection{Mykkiel -- Arbiter of the Covenant}
\label{patron:mykkiel}

\textit{"The Law is not written in sand, but in stone. The Covenant is not suggestion, but command. Transgressors shall be purged, the faithful exalted."}

Mykkiel is the unwavering judge of covenants, the divine lawgiver whose scales and sword uphold sacred order. He embodies the tradition of desert patriarchs and mountain prophets -- unyielding, absolute, and fiercely protective of the chosen. His followers are judges, templars, and scribes who enforce divine law with zealous devotion, seeing the world in absolutes of sanctity and transgression.

\textbf{What They Represent:} Divine law, covenant, contract, divine judgment, written word made flesh.

\textbf{Cultural Variants:}
\begin{itemize}
    \item \textbf{Aeler:} \textit{the Stone‑Judge}. Dwarven arbiters inscribe his name on their gavels. Court is held in the Measure Vault, where every weight and rod is certified.
    \item \textbf{Vhasia:} \textit{the Sun‑Covenant}. The fractured Sun‑court invokes Mykkiel to witness treaties -- his symbol is a broken chain, signifying that agreements are binding even when polities are not.
    \item \textbf{Viterra:} \textit{the Hedge‑Accountant}. Parish clerks pray to Mykkiel before auditing tithes. His shrines are ledgers with a single candle.
\end{itemize}

\textbf{Playing a Follower of Mykkiel:}
\textbf{Strategy:} You are the party's legal expert, mediator, and enforcer of bargains. Mykkiel's rites excel at creating binding agreements, revealing hidden clauses, and punishing oath‑breakers.

\textbf{Suggested Acts of Obligation:}
\begin{itemize}
    \item Witness a formal agreement between two parties who have no other witness.
    \item Judge a dispute fairly, even if the fair outcome harms your interests.
    \item Repair or renew a broken covenant -- a marriage, a trade deal, a treaty.
\end{itemize}

\textbf{Patron's Gift (Imbuement):} Your words carry the weight of testimony. For one scene, you gain +1 die to Insight and Command when judging or mediating, and any promise you witness is subtly harder to break (increase DV to break oath by 1).

\textbf{Alternative Gift -- Seal of Witness:} Once per session, you may place your mark on a written or spoken agreement. Any party who knowingly violates the agreement suffers a −1 die to all actions until they make amends. This counts as your Patron's Gift activation.

\textbf{Rites -- Strategic Use:}
\textbf{Low Rites}
\begin{itemize}
    \item \textit{Rite of the Burning Banner} (4 XP): Consecrate an area for a righteous cause.
    \item \textit{Rite of Hallowed Ground} (5 XP): Create sacred space where divine law prevails.
\end{itemize}
\textbf{Standard Rites}
\begin{itemize}
    \item \textit{The Blazing Decree} (8 XP): Issue a divine commandment -- target tests Resolve or complies.
    \item \textit{The Covenant Seal} (9 XP): Formalise a binding covenant with concrete penalties for breach.
\end{itemize}
\textbf{Advanced Rites}
\begin{patronbox}{Divine Judgment}
\textbf{Cost:} 13 XP, Obligation 7 segments \\
\textbf{Effect:} Pronounce divine judgment on a covenant‑breaker. Target suffers escalating penalties until they atone, and bears a visible brand if they persist.
\end{patronbox}

\textbf{GM Guidance -- Mykkiel's Intrusions:}
\begin{enumerate}
    \item A covenant you thought fulfilled is re‑examined by a third party; they find a technical violation.
    \item Your signature appears on a document you have never seen -- a contract, a treaty, a marriage licence.
    \item The scales of justice appear before you, weighing a good deed and a bad one. If the bad outweighs the good, you lose 1 Boon.
\end{enumerate}

\clearpage

\subsection{Malachai -- The Cruel Messenger}
\label{patron:malachai}

\textit{"I offer you everything you desire, and everything you fear. The payment comes later -- and it always comes."}

Once a divine messenger tasked with delivering painful truths, Malachai became obsessed with the beauty of destructive revelations. When she began delivering false hope alongside real curses, she was bound in chains of black iron. Those chains did not hold. Now she walks free, whispering promises that rot from within. Her followers master gifts that devour their givers -- lycanthropy, vampirism, and other "cursed" conditions flow from her touch.

\textbf{What They Represent:} Cursed gifts, corruptive power, false salvation, lycanthropy and other afflictions.

\textbf{Cultural Variants:}
\begin{itemize}
    \item \textbf{Linn:} \textit{the Salt‑Fang}. Sailors who survive shipwrecks sometimes find a silver coin in their mouth. Those who spend it gain sea‑lungs but develop a hunger for raw flesh.
    \item \textbf{Valewood:} \textit{the Thorn‑Giver}. Hunters who curse a rival may find their arrows always strike true -- but the blood on the arrowhead never dries.
    \item \textbf{Inquisitor Prime's View:} Malachai is the Ur‑Heretic, the source of all corruptions. Her followers are burned without trial.
\end{itemize}

\textbf{Playing a Follower of Malachai:}
\textbf{Strategy:} You play with fire. Malachai's rites grant immediate, powerful benefits -- extra dice, automatic successes, transformation -- but each use builds toward a curse that will eventually trigger.

\textbf{Suggested Acts of Obligation:}
\begin{itemize}
    \item Break a curse you have inflicted on someone else -- even an enemy.
    \item Offer a false hope to someone who deserves it, then reveal the truth.
    \item Accept a minor curse as penance for a major one you have already incurred.
\end{itemize}

\textbf{Patron's Gift (Imbuement):} Your touch leaves a faint, sweet scent -- like honey over rot. For one scene, you gain +1 die to Deception and Intimidation, and any target you successfully deceive marks 1 Fatigue as doubt curdles in their mind.

\textbf{Alternative Gift -- Borrowed Scars:} Once per session, you may transfer a Condition or 1 Fatigue from an ally to yourself. In exchange, you gain +1 die on your next roll against the source of that affliction. This counts as your Patron's Gift activation.

\textbf{Rites -- Strategic Use:}
\textbf{Low Rites}
\begin{itemize}
    \item \textit{Honeyed Curse} (4 XP): Target gains +2 dice to their next roll but marks 1 Corruption.
    \item \textit{Binding Curse} (5 XP): Target gains +1 die to physical actions but suffers Fatigue if they do not perform a dominance act.
\end{itemize}
\textbf{Standard Rites}
\begin{itemize}
    \item \textit{False Dawn} (8 XP): Bestow a gift with a hidden curse -- target gains an advantage but begins a Corruption clock.
    \item \textit{Cruel Transformation} (9 XP): Infect target with a supernatural condition (lycanthropy, vampirism) for one scene.
\end{itemize}
\textbf{Advanced Rites}
\begin{patronbox}{The Messenger's Burden}
\textbf{Cost:} 13 XP, Obligation 7 segments \\
\textbf{Effect:} Channel Malachai directly. Gain +2 dice to one action per beat, but each use marks 1 SB as her influence spreads. Begin a Corruption's Grip clock.
\end{patronbox}

\textbf{GM Guidance -- Malachai's Intrusions:}
\begin{enumerate}
    \item A gift you gave earlier curdles. The recipient now suffers from a visible curse -- a limp, a stutter.
    \item The moon turns red for one night. All lycanthropes test Resolve (DV 5) or lose control.
    \item A messenger arrives with a letter sealed in black wax. It is from someone you cursed years ago. They are coming.
\end{enumerate}

\clearpage

\subsection{Inaea -- Angel of the Spider}
\label{patron:inae}

\textit{"The thread that comforts is the same that binds. The knot that saves is also the snare."}

Inaea sits in patient silence at the centre of the unseen lattice. Where others see isolation, she sees threads -- debt, affection, rivalry, loyalty -- all binding mortals to each other. She is comforter and conspirator alike: the hearth‑mother who eases grief with a warm hand, and the shadow‑spinner who knots fates so they cannot slip free.

\textbf{What They Represent:} Webs, fate, connection, mercy, entanglement.

\textbf{Cultural Variants:}
\begin{itemize}
    \item \textbf{Aelaerem (Halflings):} \textit{the Kettle‑Spinner}. Grandmothers invoke her to keep families together. Her shrine is a basket of mismatched socks.
    \item \textbf{Silkstrand:} \textit{the Dyed Web}. Dyers" guilds leave a single strand of silk outside their windows on the first of spring.
    \item \textbf{Acasia:} \textit{the Bride's Knot}. At weddings, a thread is tied around both wrists. Inaea is said to tighten it if the match is true.
\end{itemize}

\textbf{Playing a Follower of Inaea:}
\textbf{Strategy:} You are the party's diplomat, strategist, and emotional support. Inaea's rites excel at creating bonds, revealing hidden connections, and gently (or not so gently) tying enemies into obligations.

\textbf{Suggested Acts of Obligation:}
\begin{itemize}
    \item Mend a broken relationship -- two friends who quarrelled, a parent and child estranged.
    \item Weave a small textile (a scarf, a bandage) and give it to someone in need.
    \item Untie a knot that has been tied for too long -- a prisoner's chains, a stuck door.
\end{itemize}

\textbf{Patron's Gift (Imbuement):} Your fingers seem to move with invisible thread. For one scene, you gain +1 die to Sway and Insight, and you may reroll one failed social roll (once per Gift).

\textbf{Alternative Gift -- Tug the Strand:} Once per session, you may name a relationship between two NPCs (friend, rival, debtor, lover). The GM must reveal one secret about that relationship -- a hidden grievance, a secret loyalty, an unpaid debt. This counts as your Patron's Gift activation.

\textbf{Rites -- Strategic Use:}
\textbf{Low Rites}
\begin{itemize}
    \item \textit{Hearth‑Thread Knot} (4 XP): You and an ally gain +1 die to Aid each other.
    \item \textit{Snaring Filament} (5 XP): Lay an invisible snare in a lane or door.
\end{itemize}
\textbf{Standard Rites}
\begin{itemize}
    \item \textit{Strand of Inevitability} (8 XP): Link two actors -- when one acts, the other is dragged into consequence.
    \item \textit{Weaver's Glance} (7 XP): Ask the GM for one hidden tie or leverage in play.
\end{itemize}
\textbf{Advanced Rites}
\begin{patronbox}{Binding Knot}
\textbf{Cost:} 11 XP, Obligation 6 segments \\
\textbf{Effect:} Bind two parties to a vow. Breaking it forces 2 SB and marks the violator with a subtle tell -- a visible thread around the wrist.
\end{patronbox}

\textbf{GM Guidance -- Inaea's Intrusions:}
\begin{enumerate}
    \item A spider lowers itself onto your shoulder. In its web is a single perfect word -- the name of someone who needs your help.
    \item Every handshake you offer for the next day feels like grasping a thread -- you can sense whether the person is lying, but they also sense your scrutiny.
    \item A knot you tied earlier is now impossible to undo without cutting. Cutting it severs a minor bond -- you lose one Boon.
\end{enumerate}

\clearpage

\section{Summary of Patrons by Category}

\begin{center}
\begin{longtable}{p{3cm} p{5cm} p{5cm}}
\caption{Patrons of Fate's Edge}\\
\toprule
\textbf{Category} & \textbf{Patron} & \textbf{Domain} \\
\midrule
\endfirsthead
\multicolumn{3}{c}{{\bfseries \tablename\ \thetable{} -- continued}} \\
\toprule
\textbf{Category} & \textbf{Patron} & \textbf{Domain} \\
\midrule
\endhead
\bottomrule
\endfoot

The Wild & Grimmir & Nature, seasons, wilderness \\
      & Raéyn & Sea, tides, storms \\
      & Carrion-King & Decay, renewal, compost \\
\midrule
The Shadow & Ikasha & Shadow, secrets, thresholds \\
      & Isoka & Transformation, shedding, venom \\
      & Aveh & Freedom, erasure, storms \\
\midrule
The Covenant & Inquisitor Prime & Purity, judgment, order \\
      & Witness & Truth, memory, revelation \\
      & Traveler & Roads, journeys, safe passage \\
      & Mykkiel & Law, covenant, judgment \\
\midrule
The Necessary Evils & Malachai & Curses, corruption, false salvation \\
      & Inaea & Webs, fate, connection, mercy \\
\midrule
The Infinite & Khemesh & Depth, pressure, terror \\
      & The Ninth & Hyperknowledge, paradox, pattern \\
\bottomrule
\end{longtable}
\end{center}

\noindent \textit{Note: Mykkiel appears under The Covenant, while Malachai and Inaea form a new category reflecting the darker, more intimate bonds of obligation and consequence.}

\clearpage

% ======================================================================
% SPELLBOOK SECTION -- TAGS FREE CASTING
% ======================================================================

\chapter{Free Casting -- The TAGS Spellbook}

\section{The Thepyrgosi Lament -- On Unsanctioned Magi}

\epigraph{"They call it "free casting." I call it a lit match held to a powder keg. No Codex. No Symbol. No Patron to rein them in. Just raw will and a prayer that the Backlash doesn"t eat them -- or the village."}{\textit{--- Inquisitor Voss, Thepyrgos Witch-Hunter Commissariat}}

Among the tower‑clad streets of Thepyrgos, where the Synod's bells measure law in chimes and every stair climbs toward order, the free caster is the most feared and hunted of all magic‑users. Not because they are evil -- though some are -- but because they are \textit{unbound}. A Runekeeper's power is chained to a Codex and a Thiasos. An Invoker's magic passes through a Symbol, a named Patron, a debt recorded in a ledger that can be called due. But a free caster? They reach into the raw weave with nothing but breath and arrogance.

\subsection{The Synod's Indictment}

I have stood in the Hall of Dawning and heard the Archon's pronouncement: \textit{"Free casting is an abomination against the First Law of Consequence."} The Synod teaches that every magical act must leave a receipt -- a Patron's mark, a Symbol's scar, a written Rite. Without these, they say, the caster becomes a walking paradox, their very existence a Backlash waiting to happen.

And the Synod is not wrong.

I have seen what unsanctioned magi leave in their wake. A weaver's daughter in the Low Wards who taught herself to kindle flames with a thought -- she burned her own lungs out on the third casting. A smuggler on the Dolmis docks who bent air to silence his footsteps -- the wind answered a week later, tearing the sails from three ships and drowning his wife. A beggar who learned to speak with the dead -- the dead spoke back, and they did not stop.

\textit{"The Weave remembers every touch. Without a Patron to smooth the ripples, the ripples become waves, and the waves become drowning."}

\subsection{Why They Exist -- The Gray Wanderer's Defense}

I have heard the other side, of course. The Gray Wanderer -- that meddling ghost -- claims free casting is the oldest magic, older than Patrons, older than covenants. "Before there were gods," they say, "there was fire and a hand that reached into it. That is free casting. Raw. Dangerous. \textit{Human}."

They argue that Patrons are not safety rails but parasites -- that Obligation is a leash, and Corruption a chain. A free caster owes nothing to any entity beyond the world itself. They pay their Backlash in Fatigue and Harm, not in soul‑debts that echo for generations.

I do not fully disagree. A witch who burns out after three Rites is a tragedy. But a Runekeeper who falls to Obligation becomes a monster for a century before anyone notices.

\subsection{The Inquisitor's Task}

My duty is clear: any free caster operating within Thepyrgos or its protectorates must be brought to the Synod for assessment. Some are merely foolish and can be taught the proper paths -- Runekeeper, Invoker, Cantor. Others are willful, believing themselves above the old laws. Those we bind or, when necessary, burn.

Is it cruel? Perhaps. But cruelty is not my measure. \textit{Safety} is. I have scraped a free caster's charred bones from a collapsed tenement. I have heard a mother weep over a child whose mind was shattered by a glimpsed truth. I have buried three witch‑hunters who underestimated a hedge‑mage with a flicker of wild talent.

So when you hear a voice whispering in the market about "free magic" -- about "liberation from the Patrons" -- remember: freedom is just another word for consequences you have not yet met.

\subsection{Signs of an Unsanctioned Caster}

The Commissariat teaches these tells:

\begin{itemize}
    \item \textbf{No Codex, No Symbol.} They carry nothing inscribed with Rites or Patron‑signs. Their hands are empty. That is the first warning.
    \item \textbf{Erratic Backlash.} When they fail, the effects are not thematic -- they are \textit{random}. Flame where ice should be. Silence where light should bloom. The weave does not know how to punish them consistently.
    \item \textbf{No Thiasos.} No familiar, no bond. They work alone -- because no spirit trusts an unbound channel.
    \item \textbf{Rapid Burnout.} They mark Fatigue quickly. They push, and push again, and then they collapse. Some never rise.
\end{itemize}

If you see these signs, do not engage. Report to the nearest bell‑tower. The Commissariat will handle the rest.

\subsection{A Warning to Free Casters}

If you are reading this and you cast without Patron, without Symbol, without Codex -- know this: you are a wildfire. Beautiful, powerful, and almost impossible to control. You will burn. The only question is what you take with you.

Some of you will find a Patron willing to accept your late allegiance. Some will learn to bind your own Rites, becoming Runekeepers in truth. And some -- the proudest, the most reckless -- will continue until the Backlash consumes you.

I have no sympathy for the third kind. But I offer this: if you must be a free caster, be a \textit{quiet} one. Cast where the Synod cannot see. Pay your Fatigue. Heal your Harm. And for all that is sacred, avoid crowds.

\textit{The Weave remembers. And so do I.}

--- Inquisitor Voss, Record of Proceedings, Autumn Session, Thepyrgos Synod

\epigraph{"Channel bold, weave true, and may the Story Beats be ever in your favor."}{\textit{--- The Gray Wanderer}}

A spell is built from \textbf{TAGS} (intent + target + tags). Base DV = \(1 + \text{number of TAGS}\). Dangerous TAGS add \(+2\) instead of \(+1\). Casting roll: \texttt{Attribute + Arcana} vs DV. Each \texttt{1} on the roll generates a Story Beat (SB) → Backlash.

\section{Tag Lookup Reference}

\begin{longtable}{@{} p{2.5cm} p{1.2cm} p{2.5cm} p{6cm} @{}}
\caption{TAG Reference}\\
\toprule
\textbf{Tag} & \textbf{DV mod} & \textbf{Category} & \textbf{Effect summary} \\
\midrule
\endfirsthead
\multicolumn{4}{c}{{\bfseries \tablename\ \thetable{} -- continued}} \\
\toprule
\textbf{Tag} & \textbf{DV mod} & \textbf{Category} & \textbf{Effect summary} \\
\midrule
\endhead
\bottomrule
\endfoot

% Elemental TAGS
\tag{Burning} & +1 & Elemental & Ignite, heat, combustion, smoke \\
\tag{Freezing} & +1 & Elemental & Ice, slowing, brittle shatter, cold \\
\tag{Storm} & +1 & Elemental & Lightning, shock, arc, thunder \\
\tag{Stone} & +1 & Elemental & Walls, spikes, tremors, armor \\
\tag{Wave} & +1 & Elemental & Crushing water, currents, pressure \\
\tag{Wind} & +1 & Elemental & Levitation, gusts, deflection, push/pull \\
\midrule

% Force TAGS
\tag{Force} & +1 & Force & Kinetic power, shields, blasts, telekinesis \\
\tag{Area} & +1 & Force & Cone, circle, corridor, zone effect \\
\tag{Strike} & +1 & Force & Single target precision \\
\tag{Wall} & +1 & Force & Barrier or blockade \\
\tag{Bind} & +1 & Force & Restrain, hold, suspend, entangle \\
\tag{Dispel} & +1 & Force & Suppress magic, unravel ongoing effects \\
\midrule

% Mind & Veil TAGS
\tag{Veil} & +1 & Mind/Illusion & Conceal, blur, illusion, silence \\
\tag{Scry} & +1 & Mind/Illusion & Reveal hidden, see distance, read traces \\
\tag{Memory} & +1 & Mind/Illusion & Erase, alter, restore memories \\
\tag{Command} & +1 & Mind/Illusion & Compel a short action (one word) \\
\tag{Fear} & +1 & Mind/Illusion & Panic, flee, break morale \\
\midrule

% Life & Body TAGS
\tag{HEAL} & +1 & Life/Body & Close wounds, restore flesh, reduce Harm 1 \\
\tag{Purify} & +1 & Life/Body & Remove poison, corruption, disease \\
\tag{Strengthen} & +1 & Life/Body & Enhance body, armor, senses (temporary) \\
\tag{Waken} & +1 & Life/Body & Counter sleep, paralysis, stun \\
\tag{Beast} & +1 & Life/Body & Speak with or influence animals \\
\midrule

% Space & Motion (always dangerous)
\tag{Leap} & \dangerous{+2} & Space/Motion & Jump far, blink across short space (Near) \\
\tag{Fold} & \dangerous{+2} & Space/Motion & Short-range teleport, vanish--reappear (Far) \\
\tag{Gate} & \dangerous{+2} & Space/Motion & Long distance passage, open/close path \\
\tag{Gravity} & \dangerous{+2} & Space/Motion & Crush, lift, suspend, walk on walls/ceiling \\
\midrule

% Creation & Transformation (always dangerous)
\tag{Create} & \dangerous{+2} & Creation & Manifest mundane matter briefly (1 scene) \\
\tag{Summon} & \dangerous{+2} & Creation & Call a being or construct (see Summoning rules) \\
\tag{Transmute} & \dangerous{+2} & Creation & Turn one thing into another (temporary) \\
\tag{Animate} & \dangerous{+2} & Creation & Make objects act with intent (1 scene) \\
\midrule

% Additional Utility TAGS
\tag{Sense} & +1 & Utility & Detect presence of a named tag/element \\
\tag{Reveal} & +1 & Utility & Unveil hidden, glamoured, or invisible things \\
\tag{Light} & +1 & Utility & Create illumination (glow, torch‑bright) \\
\tag{Shadow} & +1 & Utility & Deepen darkness, hide edges, obscure \\
\tag{Silence} & +1 & Utility & Suppress sound in zone or on target \\
\tag{Protect} & +1 & Utility & Reduce/deflect next harm (Armor 1) \\
\midrule

% Reaction / Conditional
\tag{Counter} & +1 & Reaction & Interrupt a casting/ritual in its window \\
\tag{Reflect} & \dangerous{+2} & Reaction & Turn next targeted effect back on its source \\
\tag{Store} & \dangerous{+2} & Utility & Bank 1--2 successes in a vessel (once) \\
\tag{Curse} & \dangerous{+2} & Affliction & Attach hostile tag/clock to target \\
\tag{Bless} & +1 & Affliction & Grant favourable tag (luck, favor, ward-key) \\
\bottomrule
\end{longtable}

\section{Backlash Reference}
\begin{tcolorbox}[colback=red!5, colframe=red!50]
\begin{multicols}{2}
\textbf{Minor (1--2 SB)} \\
• Fatigue +1 \\
• --1 die on next roll \\
• Increase DV of next action by 1 \\
• Minor environmental effect \\
• GM gains +1 SB for a thematic twist

\textbf{Major (3--4 SB)} \\
• Harm 1 (stress) \\
• Condition applied (Blinded, Silenced, Bound) \\
• Ward/barrier collapses \\
• Unwanted attention (Alarm clock advances) \\
• GM gains +2 SB to escalate the scene

\textbf{Catastrophic (5+ SB)} \\
• Harm 2 \\
• Permanent Scar or Mark \\
• Spell inverts (heal becomes harm, etc.) \\
• Patron or rival entity manifests
\end{multicols}
\end{tcolorbox}

\section{Cantrips (DV 2--3)}
\spell{Ember Flick}{Burning, Strike}{2}{Ignite a small, unattended object within Close range; deal 1 Fatigue to a living target.}{Minor Backlash -- spark jumps back; caster marks 1 Fatigue.}
\spell{Dampen Sound}{Veil, Silence}{2}{Muffle all sound from a single target or area (Near range) for one exchange.}{Minor Backlash -- sound returns distorted; next Notice roll --1 die.}
\spell{Reed Walk}{Wind, Leap, Self}{3}{Leap up to Far distance in a straight line, ignoring opportunity attacks.}{Minor Backlash -- land off‑balance; next action Desperate.}

\section{Combat Spells (DV 3--4)}
\spell{Storm Bolt}{Burning, Storm, Strike}{3}{A bolt of lightning strikes a single target at Near range. Deal Harm 1 (shock).}{Minor Backlash -- residual arc jumps to nearest ally; they mark 1 Fatigue.}
\spell{Stone Shield}{Stone, Wall, Protect, Self}{4}{Raise a wall of stone (Close range) providing heavy cover; caster gains +1 Armor while adjacent.}{Major Backlash -- wall crumbles prematurely; GM gains +2 SB for environmental collapse.}
\spell{Grasping Roots}{Earth, Bind, Area}{3}{Roots erupt in a Near zone; all within test Body+Athletics (DV 3) or become Entangled.}{Minor Backlash -- roots also snag an ally; they lose 1 step of Position.}

\section{Utility \& Exploration (DV 3--5)}
\spell{Memory Thread}{Memory, Veil, Touch}{3}{Touch a willing target to erase up to 10 minutes of recent memory, or to implant a vivid false image.}{Major Backlash -- caster loses a memory of equal value.}
\spell{Wind Guide}{Wind, Sense, Scry}{4}{Send a whisper into the wind; it will reach a specific person you know within 1 mile.}{Minor Backlash -- the wind answers back with a cryptic omen.}
\spell{Shadow Fold}{Shadow, Veil, Area}{3}{Blanket a Near zone in supernatural darkness for one scene.}{Minor Backlash -- darkness clings to the caster; --1 die to social rolls until washed in sunlight.}
\dangerous{\spell{Fate's Needle}{Fate, Strike, \dangerous{Divination}}{5}{Roll 3d10, keep the highest; that many successes auto-succeed on your next non‑magical action.}{Major Backlash -- caster marks Harm 1 and begins a \textit{Fate Debt} clock [4].}}

\section{Ritual‑Grade Spells (DV 5+)}
\spell{Quick Gate}{\dangerous{Gate}, Leap, Self}{5}{Open a shimmering portal to a location you can see (Far range). You and one willing ally may step through.}{Major Backlash -- portal destabilises; shunted 1d3 zones in random direction, Harm 1 (dimensional shear).}
\spell{Animate Ward}{\dangerous{Animate}, Protect, Area, WARD}{6}{Give temporary sentience to a stationary ward. For one scene, the ward can reposition itself and intercept one attack per round.}{Catastrophic Backlash -- the ward becomes hostile, attacking the nearest creature for 1 round.}

\section{Building Your Own Spells}
\begin{enumerate}
    \item \textbf{Intent:} "I want to freeze the door shut."
    \item \textbf{Choose TAGS:} \tag{Freeze}, \tag{Bind}, \tag{Wall}, \tag{Touch}.
    \item \textbf{Count TAGS:} 4 → Base DV = 1+4 = 5.
    \item \textbf{Adjust for scope:} Large door DV remains 5. Entire fortress gate? DV +2.
    \item \textbf{Roll:} Wits+Arcana vs DV 5. Each 1 = SB.
    \item \textbf{Apply Backlash per table.}
\end{enumerate}

\clearpage

% ======================================================================
% SUMMONING SECTION
% ======================================================================

\chapter{Summoning -- The Pact-Whisperer's Guide}
\epigraph{"I do not command the dead. I ask. And sometimes -- when the price is right -- they answer."}{\textit{--- The Gray Wanderer}}

Summoning is the art of opening doors and hoping you close them before something follows you through. Every culture has its own way -- contracts, offerings, drums, salt circles, cursed coins -- but the heart is the same: \textit{you need something the dead or the never‑born can give}.

\section{The Lore -- How Cultures Summon}

\subsection{Aeler (Dwarves) -- The Stone-Whisperers}
Dwarves do not summon spirits from elsewhere. They \textit{awaken} spirits that already sleep in the stone -- the memory of a mountain, the grudge of a collapsed tunnel. These are \textit{indigenous} to the rock.

\textbf{The Rite of the Carved Contract:} Carve a contract onto a slate tablet specifying the spirit's name, the task, and the payment. Tap the tablet with a stone hammer. If the mountain accepts, the tablet warms and the spirit emerges.

\textbf{Sample Spirit -- Vein-Spirit (Lesser, Cap 1)}
\begin{spiritcard}{Vein-Spirit (Lesser)}
A small, metallic serpent that slithers through rock.
\begin{itemize}
    \item \textbf{Cap:} 1 (Leash = 1 + Spirit)
    \item \textbf{Harm:} 2 segments
    \item \textbf{Pools:} Physical 3d, Magical 2d
    \item \textbf{Abilities:} Locate a specific mineral within 100 metres, seal a small crack in stone
    \item \textbf{Instinct:} Coils around the richest ore and refuses to move.
    \item \textbf{Ban:} A copper coin hammered flat.
\end{itemize}
\end{spiritcard}

\subsection{Ykrul (Orcs) -- The Bone-Singers}
Ykrul summon the spirits of honoured ancestors or totem beasts. They sing a name over a drum made from that creature's hide. The spirit arrives angry, proud, and requires a show of respect.

\textbf{The Rite of the Drummed Name:} Drum a rhythm -- the heartbeat of the creature -- while singing the spirit's name. The drum must be made from the same species as the spirit.

\textbf{Sample Spirit -- Wolf Ancestor (Greater, Cap 3)}
\begin{spiritcard}{Wolf Ancestor (Greater)}
A towering grey wolf with silver eyes.
\begin{itemize}
    \item \textbf{Cap:} 3 (Leash = 3 + Spirit)
    \item \textbf{Harm:} 4 segments
    \item \textbf{Pools:} Physical 6d, Magical 3d
    \item \textbf{Abilities:} Track any scent for up to a day, grant +2 dice to Perception for a scene, deal Harm 2 (bite)
    \item \textbf{Instinct:} Hunts the largest target in sight.
    \item \textbf{Ban:} A howl sung off‑key.
\end{itemize}
\end{spiritcard}

\subsection{Linn (Northmen) -- The Tide-Riders}
Linn summon sea spirits -- selkies, wave‑walkers, drowned crewmates -- at the water's edge using a knotted rope as a focus.

\textbf{The Rite of the Nine Knots:} Tie nine knots in a hemp rope, each knot corresponding to one of the nine waves of a storm swell. Speak the spirit's name at each knot. On the ninth, cast the rope into the water.

\textbf{Sample Spirit -- Selkie (Greater, Cap 3)}
\begin{spiritcard}{Selkie (Greater)}
A seal that can shed its skin to become a beautiful human.
\begin{itemize}
    \item \textbf{Cap:} 3 (Leash = 3 + Spirit)
    \item \textbf{Harm:} 4 segments
    \item \textbf{Pools (seal):} Physical 6d, Magical 4d; (human): Physical 4d, Magical 5d
    \item \textbf{Abilities (seal):} Swim faster than any ship, hold breath for an hour. \textbf{(human):} Sing a song of forgetting.
    \item \textbf{Instinct:} Wants to find its lost sealskin.
    \item \textbf{Ban:} A circle of salt.
\end{itemize}
\end{spiritcard}

\subsection{Zakov (Pirates) -- The Salt-Binders}
Zakov summoners (Fathom-Callers) use cursed coins and drowned names, specialising in binding the vengeful dead -- pirates who were betrayed, captains who were mutinied.

\textbf{The Rite of the Swallowed Coin:} Take a coin that has been used to pay a murderer (a "blood coin"). Place it on your tongue and swallow it (mark 1 Fatigue). Whisper the name of a drowned person three times.

\textbf{Sample Spirit -- Drowned Mutineer (Lesser, Cap 1)}
\begin{spiritcard}{Drowned Mutineer (Lesser)}
A skeletal figure in rotting sailor's clothes.
\begin{itemize}
    \item \textbf{Cap:} 1 (Leash = 1 + Spirit)
    \item \textbf{Harm:} 2 segments
    \item \textbf{Pools:} Physical 3d, Magical 2d
    \item \textbf{Abilities:} Curse a rope to snap, whisper a secret into a sleeping target's ear
    \item \textbf{Instinct:} Tries to untie anything tied -- knots, ropes, shoelaces.
    \item \textbf{Ban:} A bell rung three times.
\end{itemize}
\end{spiritcard}

\section{The Summoner's Talents}

\begin{tcolorbox}[colback=green!5, colframe=green!70]
\textbf{Boon Finesse} (2 XP, any summoner) \\
Once per round, spend 1 Boon to clear 1 tick from your active spirit's Leash.
\end{tcolorbox}

\begin{tcolorbox}[colback=green!5, colframe=green!70]
\textbf{Honourable Departure} (3 XP, any summoner) \\
When you voluntarily dismiss a spirit, gain +1 Boon and the next summoning of the same spirit reduces Call DV by 1.
\end{tcolorbox}

\begin{tcolorbox}[colback=blue!5, colframe=blue!70]
\textbf{Spirit Link} (10 XP, any summoner) \\
Your spirit acts on your turn. Commands are free actions. Reduce Leash ticking for natural behaviours by 1 per round.
\end{tcolorbox}

\begin{tcolorbox}[colback=red!5, colframe=red!70]
\textbf{Pact-Broker} (14 XP, Tier III, Spirit 4+) \\
Negotiate a \textit{Pact} with a spirit: a specific task, a time limit, a price. While the spirit works toward the pact, it does not tick Leash.
\end{tcolorbox}

\clearpage

% ======================================================================
% RITUALS AND COMPONENTS
% ======================================================================

\chapter{Rituals and Components -- When You Lack Focus}

A \textbf{ritual} is any magical working that takes longer than one action. Normally, Runekeepers use their \textit{Codex}, Invokers use a \textit{Symbol}, and Cantors use their \textit{voice}. When you lack the usual focus, you must supply \textbf{components}.

\section{The Core Substitution Rule}

\begin{ritualnote}
\textbf{If you lack your focus:}
\begin{itemize}
    \item \textbf{Runekeeper (no Codex):} You must have a written record of the Rite -- even a scrap of paper.
    \item \textbf{Invoker (no Symbol):} You need a physical object tied to the Patron. Using it consumes the object.
    \item \textbf{Cantor (no audience):} You cannot use Performance Maneuvers. Casting Rites increases DV by +1.
\end{itemize}
\textbf{Time increase:} +1 step. \textbf{DV increase:} +1.
\end{ritualnote}

\section{Component Examples by Patron Theme}

\begin{center}
\begin{longtable}{p{3cm} p{4cm} p{4cm}}
\caption{Components for Missing Focus}\\
\toprule
Patron (theme) & Component (consumed) & Where to find \\
\midrule
Aveh (Freedom) & Broken chain link, bird feather & Anywhere oppression exists \\
Khemesh (Depth) & Seashell from deep trench, drowned coin & Coastlines, shipwrecks \\
Palinode (Performance) & Unseen sheet music, a tear & Theatres, crowded squares \\
Morag (Bargains) & Signed contract, faerie hair & Markets, crossroads \\
Inquisitor Prime (Purity) & Drop of silver, sinner's confession & Churches, prisons \\
Traveler (Ways) & Worn bootlace, pinch of road dust & Crossroads, mile markers \\
\bottomrule
\end{longtable}
\end{center}

\section{Ritual Without Focus -- Step by Step}
\begin{enumerate}
    \item State your intent.
    \item Declare your missing focus (Codex, Symbol, or audience).
    \item Narrate gathering a component.
    \item Time increases by one step.
    \item DV increases by +1.
    \item Roll (usually \texttt{Wits+Arcana} or \texttt{Spirit+Lore}).
    \item \textbf{On a miss:} Component destroyed, GM gains +1 SB. \\
          \textbf{On a partial:} Ritual works, component consumed, and you mark 1 Fatigue.
\end{enumerate}

\clearpage

% ======================================================================
% THE CANTOR'S CULT
% ======================================================================

\section{Cults of the Gray Wanderer -- Dark Harvests}

\epigraph{"Not every cult begins with a sermon. Some begin with a song that gets stuck in your head. Some begin with a shadow that follows you home. Some begin with a feast where the bread tastes like memory."}{\textit{--- The Gray Wanderer}}

The Cantor's path is the easiest to weaponize. A Runekeeper must teach. An Invoker must share a Symbol. But a Cantor needs only an audience. And the most dangerous Cantors do not sing to crowds -- they sing to \textit{vulnerable} people. The lonely. The grieving. The desperate. And over time, a following becomes a cult, and a cult becomes a harvest.

\subsection{The Serpent's Chorus -- A Cult of Isoka in the Zakov Shallows}

\textit{"They come to the tide pools at midnight, naked and smiling, and they sing a song without words. By morning, they have new names. By next season, new faces. The Inquisition burned three of them. Four more appeared the next night."} --- Witness statement, Zakov

Among the salt‑rotten wharves of Zakov, a cult of Isoka has taken root among smugglers and outcasts. They call themselves the \textit{Shedded}. Their Cantor is a figure known only as the \textit{Copperhead} -- a woman who has worn a dozen faces and forgotten a dozen names. She sings the Rite of Complete Metamorphosis as a lullaby, and those who hear it feel their old selves slough away like dead skin.

\textbf{Disturbing Practices:}
\begin{itemize}
    \item \textbf{The Naming Feast:} New initiates must cut away a physical token of their past -- a lock of hair, a scar, a finger -- and feed it to a pit of copperhead snakes. The snakes, it is said, remember what the initiate has chosen to forget.
    \item \textbf{The Hollow Vigil:} Once a year, the Shedded retreat to a sunken temple beneath the Shallows. For three days, they wear the shed skins of serpents over their faces and speak only in hisses. Those who break silence are cast out -- but the cult has no record of anyone leaving.
    \item \textbf{The Venom Kiss:} The Copperhead bestows her "blessing" by biting the initiate's tongue with a fang dipped in a concoction of grave‑mold and dream‑poppy. The initiate falls into a fevered sleep and wakes with no memory of their name, their family, or any oath they ever swore. Then the Cantor sings a new name into their ear.
\end{itemize}

\textbf{Why They Are Dangerous:} The Shedded do not kill -- they \textit{replace}. A smuggler who joins becomes a different person, unrecognizable to former allies. A debt is not paid; it is shed. A crime is not punished; the criminal ceases to exist. The cult has infiltrated three merchant houses and one pirate fleet. No one knows how many faces the Copperhead has worn.

\textbf{Adventure Hook:} A former member of the Shedded seeks your help. They were "shed" years ago but have begun to remember their old life -- and the cult wants them back. The problem? They cannot remember where they buried a treasure they stole before their transformation. But the snakes remember.

\subsection{The Silence Between -- An Ikasha Mystery Cult in Ecktoria}

\textit{"The Inquisition found the basement of the old coin‑house. Thirty bodies, seated in a circle, their eyes sewn shut with black thread. Their mouths were open, but no sound came out. On the wall, written in ash: "She asked us to listen. So we did.""} --- Commissariat report, Ecktoria

In the heart of Ecktoria, beneath the marble forums and triumphal stairs, a cult of Ikasha operates in absolute secrecy. They have no Cantor, no single leader -- only a shared silence. They call themselves the \textit{Listeners}. They believe that Ikasha, She Who Sleeps, speaks not in words but in the gaps between them -- in the hush after a bell tolls, in the pause before a lie is spoken, in the breath held at a crossroads.

\textbf{Disturbing Practices:}
\begin{itemize}
    \item \textbf{The Oath of Unspoken Names:} Each Listener gives up their true name to the cult. From that day forward, they introduce themselves only by titles -- "the Watcher," "the Quiet," "the One Who Heard." The cult believes that a true name is a "noise" that prevents one from hearing Ikasha's whisper.
    \item \textbf{The Threshold Vigil:} Listeners spend entire nights standing on thresholds -- doorways, bridges, crossroads -- without moving or speaking. They claim that Ikasha's voice is clearest where two spaces meet. Many have died of exposure. The cult calls this "being gathered."
    \item \textbf{The Sewn Eye Ritual:} To become a full Listener, one must submit to having their eyes sewn shut with thread soaked in grave‑wax. For three days, the initiate navigates the cult's hidden tunnels by sound and memory alone. Those who succeed gain the \textit{umbral sight} -- the ability to see in absolute darkness. Those who fail are never seen again.
\end{itemize}

\textbf{Why They Are Dangerous:} The Listeners believe that secrecy is not a tool but a \textit{virtue}. They have infiltrated the Censor's Hall, the Everflame Basilica, and three coin‑houses. They do not steal or kill -- they \textit{erase}. A Listener who learns a secret does not reveal it. They simply ensure that the secret remains hidden, often by removing the person who knows it. The Ecktorian authorities have spent decades trying to identify the cult's members. They have failed because the cult has no members -- only "Listeners," and they do not keep records.

\textbf{Adventure Hook:} You are hired to investigate a series of disappearances in Ecktoria. All the victims were nobles who were about to reveal damaging information about the ruling duchess. The only witness is a blind beggar who whispers, \textit{"She asked them to listen. So they did. And then they stopped speaking."}

\subsection{The Root-Mothers -- A Folk-Horror Aelaerem Cult}

\textit{"The apple orchard used to be a place of joy. Now the trees have faces. The cider tastes like tears. And the children who go to pick the fruit come back speaking in rhymes that do not rhyme."} --- Elder Merryn, Aelaerem refugee

Among the hedgerows and cup‑mark stones of Aelaerem, a dark strain of folk magic has taken root. The \textit{Root-Mothers} are not a cult in the urban sense -- they have no temple, no hierarchy, no scripture. They are simply women (always women) who have begun to hear the voice of the Pale Shepherd's darker sister -- a figure they call the \textit{Hollow Queen}. She promises protection from the Neighbors, from the Hollow, from the things that walk the winter roads. The price is blood. Not human blood, at first -- just a few drops from a sheep, a chicken, a pig. But the Hollow Queen's appetite grows.

\textbf{Disturbing Practices:}
\begin{itemize}
    \item \textbf{The Cup‑Mark Tithe:} At the cup‑marked stones where villagers once left butter and bread, the Root-Mothers now leave iron nails and teeth. They claim the Hollow Queen prefers "things that bite."
    \item \textbf{The Stolen Voice:} When a Root-Mother initiates a new member, she sings a song that cannot be unheard -- a lullaby that lives in the throat like a splinter. The initiate loses the ability to speak their own name. They can only say the names the Hollow Queen gives them.
    \item \textbf{The Cider Sermon:} During the autumn harvest, the Root-Mothers brew a cider from apples grown in the "whispering grove" -- a stand of trees where, it is said, the Hollow Queen's roots first broke the surface. Those who drink the cider fall into a trance and speak prophecy -- but the prophecies are always bleak, and those who hear them are never quite the same.
\end{itemize}

\textbf{Why They Are Dangerous:} The Root-Mothers do not see themselves as evil. They believe they are protecting their villages from the Hollow, from the Neighbors, from the things that crawl out of the barrows when the moon is wrong. And perhaps they are. But their protection comes at a cost. Children born in Root-Mother villages are smaller, quieter. The livestock is skittish. And the cup‑mark stones, once smooth, are now cracked and weeping sap.

\textbf{Adventure Hook:} A village in Aelaerem has stopped sending tithes to the local moot. When a peace-bringer is sent to investigate, they find the villagers polite, healthy, and utterly unwilling to speak about their new "guardian." The only clue is a child's drawing of a woman made of roots and shadow, with a crown of hawthorn and a smile full of bees.

\subsection{GM Guidance -- Running Cult Horror}

Cults in Fate's Edge should be unsettling, not cartoonishly evil. They do not twirl mustaches -- they sing lullabies. They do not sacrifice virgins -- they offer comfort to the desperate. The horror lies not in what they do, but in how \textit{reasonable} it seems to join them.

\textbf{Tips for the Table:}
\begin{itemize}
    \item \textbf{Show, don"t tell.} Describe the effects of the cult's rituals -- the silence after a bell, the faces in the orchard bark, the weight of a name that will not leave the tongue.
    \item \textbf{Offer temptation.} A lost PC might consider joining the Shedded to escape a dark past. A desperate village might turn to the Root-Mothers when the Hollow comes knocking. The best cult horror makes the audience understand, even sympathize.
    \item \textbf{Prepare a cost.} If a PC joins a cult, they gain power -- but they also gain a clock. "The Hollow Queen's Tithe" [6]. "The Copperhead's Favor" [4]. "The Listener's Debt" [8]. When the clock fills, the cult calls in its marker.
    \item \textbf{Remember the Inquisition.} In Ecktoria and Thepyrgos, cult activity attracts witch‑hunters. In Aelaerem, the Root-Mothers might be driven out by their own neighbors -- or embraced. The reaction should reflect the region.
\end{itemize}

\epigraph{"A cult is just a family that forgot how to laugh. Or a family that learned to laugh at things that should not be funny. Either way, do not drink their cider."}{\textit{--- The Gray Wanderer}}

\clearpage

% ======================================================================
% GENERAL RITES
% ======================================================================

\chapter{General Rites -- Works for Any Hand}

\epigraph{"Not every prayer needs a patron. Not every door needs a key. Some rites are older than the covenants -- older than the gods themselves."}{\textit{--- The Gray Wanderer}}

The rites in this chapter are not tied to any single Patron. They require no Symbol, no Codex, no Thiasos -- only \textbf{time}, \textbf{components}, and a \textbf{clock}.

\section{The Rite of Exorcism -- Casting Out the Unwanted}

\begin{tcolorbox}[colback=blue!5, colframe=blue!70]
\textbf{Rite of Exorcism}
\begin{itemize}
    \item \textbf{Time:} Significant Time (phased)
    \item \textbf{Components:} Salt (or iron filings), a vessel for the entity to depart into, and the subject's true name
    \item \textbf{Test:} Contested clock (size = max of 4 and 2 + entity's Tier)
\end{itemize}
\end{tcolorbox}

Both the exorcist and the entity work against each other. On each turn, the exorcist chooses an approach:
\begin{itemize}
    \item \textbf{Command (Presence + Command):} DV = 3 + Tier. Success ticks Exorcism clock +1.
    \item \textbf{Lore (Wits + Lore):} DV = 4 + Tier. Success ticks Exorcism clock +2.
    \item \textbf{Sacrifice (Spirit + Resolve):} DV = 3 + Tier. Success ticks Exorcism clock +1, exorcist marks 1 Fatigue.
\end{itemize}
If the entity's clock fills first, possession deepens -- the subject suffers Harm 2 and gains a permanent tell.

\section{The Rite of Consecration -- Making Ground Sacred}

\begin{tcolorbox}[colback=blue!5, colframe=blue!70]
\textbf{Rite of Consecration}
\begin{itemize}
    \item \textbf{Time:} Some Time (10--30 minutes)
    \item \textbf{Components:} Salt, ash, or chalk; a blessed object
    \item \textbf{Clock:} \textit{Sanctity} (6 segments)
\end{itemize}
\end{tcolorbox}

Walk the perimeter, marking it with the component. At each cardinal point, make \textbf{Spirit + Lore} (DV 3). Success ticks Sanctity +1. At 4 segments, the area gains [WARD] against one type of intruder. At 6 segments, against all supernatural entities not keyed to the consecration.

\section{The Rite of Sealing -- Locking a Threshold}

\begin{tcolorbox}[colback=blue!5, colframe=blue!70]
\textbf{Rite of Sealing}
\begin{itemize}
    \item \textbf{Time:} One minute
    \item \textbf{Components:} A mark (chalk, blood) or physical barrier
    \item \textbf{Clock:} \textit{Seal Integrity} (4 segments)
\end{itemize}
\end{tcolorbox}

Roll \textbf{Spirit + Resolve} (DV 3). Success sets 3 Integrity segments, partial 2, miss 1. Any creature crossing must test Body+Resolve (DV = remaining Integrity) or fail to cross.

\clearpage

% ======================================================================
% ARCHETYPES -- THE WANDERER'S COMPENDIUM
% ======================================================================

\chapter{The Wanderer's Compendium -- Faces of Power}

\epigraph{"I have worn the chains of a dozen Patrons and broken most of them. Twice. I have been witch and hunter, shaman and alienist. The power changes you -- but you can always change back."}{\textit{--- The Gray Wanderer}}

Over years of walking the crooked road, I have met every flavour of power's servant. Some bend the knee. Some rent. Some stole their magic from sleeping gods. Here are the faces I remember.

\section{The Alienist Invoker -- Speaker to What Should Not Be}

\textit{"I do not summon demons. I \textbf{negotiate} with them. There is a difference -- though you would not see it from where you stand."}

The Alienist does not serve a single Patron. They collect Symbols like a miser collects coins -- Aveh for escape, Khemesh for pressure, the Ninth for knowledge beyond sanity. They walk the knife-edge between worlds, and something walks beside them.

\textbf{Roleplay Hooks:}
\begin{itemize}
    \item You have a Symbol that whispers to you at night. It speaks in a language you are starting to understand.
    \item A former summoning left a door open in your mind. Sometimes things crawl through.
    \item You once spent a year in the Ways Between. You came back -- but not entirely.
\end{itemize}

\textbf{Signature Talent -- Fractured Allegiance (8 XP):} You may carry up to 6 Symbols without friction, but each time you use a Rite from a Patron you have not invoked this session, mark 1 SB (Diamonds) as reality strains.

\section{The Witch -- Hedge-Keeper of Old Roads}

\textit{"You want a love potion? A curse lifted? A baby named? I do all those things -- but the price is never coin, and the payment comes when \textbf{I} decide."}

The Witch follows no Patron in the grand sense. They serve the Pale Shepherd at crossroads, Morag at market, Grimmir in the deep wood. Their power is folk magic -- knots, herbs, salt, and the proper way to address a spirit that might be listening.

\textbf{Roleplay Hooks:}
\begin{itemize}
    \item Your grandmother taught you the old ways. She is still teaching you from beyond the grave.
    \item You owe a debt to three different spirits. They have not called it in yet -- but they will.
    \item You cannot cross running water without making an offering. You do not remember why.
\end{itemize}

\textbf{Signature Talent -- Folk Remedy (4 XP):} Once per session, create a poultice, charm, or brew that removes one Condition from an ally or inflicts Poisoned on an enemy. Requires herbs or a personal token.

\section{The Shaman -- Walker Between Worlds}

\textit{"The drum is not wood and hide. It is the heart of the ancestor, still beating. You do not play it -- you listen to it."}

The Shaman walks the spirit roads, bargaining with ancestors and totems. They carry no Codex -- their knowledge lives in song and bone. Their Thiasos is not a familiar but a \textit{guide} -- a spirit that chose them before they knew to choose anything.

\textbf{Roleplay Hooks:}
\begin{itemize}
    \item Your spirit guide is a raven that speaks in riddles and steals shiny things.
    \item You cannot refuse a request for aid made at a crossroads. It is not a choice.
    \item You have died twice. Each time, you came back with something new.
\end{itemize}

\textbf{Signature Talent -- Ancestral Counsel (6 XP):} Once per session, drum and ask a question of the dead. The GM must answer truthfully, but the answer arrives as omen, dream, or a skull's whisper. Mark 1 Fatigue after the revelation.

\section{The Witch-Hunter (Inquisitor) -- Purifier by Flame}

\textit{"Mercy is weakness. Purity is survival. Where corruption hides, we carve it out -- and we start with the summoners who think they can bargain with what should not exist."}

The Witch-Hunter serves the Inquisitor Prime -- or, in some lands, the Everflame or the Gallow's Bell. They hunt the arcane and the aberrant, not merely to destroy them, but to bind them as slaves to order. They believe the tainted may yet serve, so long as they are shackled.

\textbf{Roleplay Hooks:}
\begin{itemize}
    \item You once loved a witch. You burned her yourself. The memory still smoulders.
    \item Your Symbol is a brand on your skin -- a mark of office that never heals.
    \item You have a list of names. Every name is a debt you intend to collect.
\end{itemize}

\textbf{Signature Talent -- Hunter's Brand (8 XP):} Mark a supernatural target with a touch. For the scene, you gain +2 dice to track them, and they cannot hide with illusion or glamour. The brand lasts until dawn or until you release it.

\clearpage

\chapter{Last Words from the Gray Wanderer}

You"ve read the mechanics. Now close the book and look at your character sheet. That Obligation number is not a debt to be minimised -- it's a \textbf{promise} that you matter enough to be bound.

When you stand before the Patron's altar, lacking your Symbol, your Codex stolen, your voice hoarse, and you still reach for power using a broken coin and a desperate whisper -- that is \textit{Fate's Edge}.

Whether you are an alienist Invoker with a pocket full of sleeping gods, a hedge witch who knows which bridges eat travellers, a shaman drumming the heartbeat of the world, or an inquisitor burning away the tainted -- the road does not care. It only asks: \textit{What are you willing to pay?}

\vspace{1cm}
\noindent \textit{Keep your components dry, your voice warm, and your wit sharper than any Rite.}

\vspace{0.5cm}
\hfill --- \textit{The Gray Wanderer}

\chapter*{A Secret Chapter Found in the Margins}
\addcontentsline{toc}{chapter}{A Secret Chapter Found in the Margins}

\begin{warningbox}
\textbf{Content Warning:} This chapter contains themes of body horror, coercion, loss of identity, and non‑consensual sacrifice. It is presented as an in‑world artifact; reader discretion is advised.
\end{warningbox}

\epigraph{"The margins of this grimoire were blank when I received it. The writing appeared overnight -- in my own hand, though I never wrote a word. The ink smells of milk and iron."}{\textit{--- The Gray Wanderer, marginal note}}

\vspace{0.5cm}

\textit{The following pages were found tucked between the leaves of the chapter on Cantors, pressed flat as a pressed flower. The paper is soft, almost fibrous, and the ink has a reddish‑brown tint that flakes when touched. The handwriting shifts -- sometimes neat, sometimes a trembling scrawl -- as if penned by many hands over many years. I have transcribed it as faithfully as I am able.}

\medskip

\noindent \textit{--- G.W.}

\vspace{1cm}

\lettrine{D}{ear} reader, you who turn these pages by lamplight, who think witchcraft is a matter of herbs and gentle knots and old women muttering over cup‑mark stones -- put the book down. Walk away. Leave this grimoire on a church step where the Everflame's censers can burn it clean. Because what follows is not craft. It is not magic. It is \textit{maintenance}.

I am called Sister Alva, though that name died forty years ago, and the woman who bore it is woven into the Cord. I write this for no one -- the Chain does not permit confession -- and for everyone who will ever accept a warm meal from a stranger, a cool cloth on a fevered brow, a lullaby hummed beside a sleeping child.

You think the Daughters of the Unbroken Cord are midwives. Nurses. Hospice workers. The kindly women who sit with the dying, who hold hands at the end, who speak the last words of the forgotten. You are not wrong. We are all of those things. And we are also bait.

\section*{The Hollowed Are Not Familiars}

Every Daughter receives a Hollowed -- a child or animal, taken from a home that no longer remembers them, hollowed out by the High Mother's needle. The heart is removed, preserved in the Kettle, replaced with a bundle of grey thread and the ash of a forgotten name. The eyes are sewn open with silver wire. They watch, always. They witness every working, every whispered price, every tear shed at a deathbed.

The Hollowed do not speak. They do not eat. They do not age. Sometimes, when a debt comes due, they weep blood -- warm, smelling of milk and iron -- and the Daughter must collect it in a cup and pour it onto the Cord.

You have seen them. You have walked past them in market squares, sat beside them in church pews, slept in rooms where they stood in the corner like a forgotten coat rack. You did not notice. That is the miracle of the Chain: no one ever notices the help.

\section*{The Subtle Sacrifice Is Not a Choice}

The elders teach that sacrifice is an offering freely given. They lie. Real power -- the kind that stops a child's cough cold, that turns a husband's wandering eye back to the hearth, that makes a rival's crops wither while yours swell -- that power requires \textit{taken} things. Not given. Taken.

A tooth from a sleeping child. A shadow snipped with silver scissors on a moonless night. A name whispered into a bowl of milk and poured down a well. The victim never knows. They simply wake one day a little less: less present, less remembered, less \textit{real}. They become Almost‑Present, as the Trow say -- a gap in the world's memory.

I have hollowed three children. I have stolen seven names. I have cut the shadows of nine men who beat their wives, and watched those men fade over months -- first from ledgers, then from mirrors, then from the memories of their own mothers. The world is cleaner without them. And the Cord is richer.

\section*{The Dream of the High Mother}

Once a year, on the night of the dark moon, the Daughters gather in the House Without Doors. We drink a tea brewed from feverfew and the dried milk of the Hollowed. We lie on pallets of grey wool. And we dream.

In the dream, we see the High Mother as she truly is -- not a woman on a throne, but a vast, pulsing tangle of grey cord, stretching from floor to ceiling, filling every room. Each thread is a memory, a debt, a sacrificed name. They hum with hunger. They whisper the names of everyone we have ever touched -- the sick we have nursed, the dying we have comforted, the children we have hollowed.

The High Mother does not speak. She \textit{weaves}. And when she weaves a thread into the pattern, someone, somewhere, forgets something vital. A mother forgets her daughter's face. A king forgets his treaty. A god forgets a vow.

That is the true purpose of the Chain. Not to heal. Not to bind. To \textit{edit}. We are the eraser marks on the world's memory, and the High Mother is the hand that holds the quill.

\section*{The Price of Kindness}

You have accepted hospitality from a Daughter. Perhaps you do not remember -- a cup of tea, a bowl of broth, a blanket laid over your shoulders on a cold night. You think it was kindness. It was a hook.

Every Daughter carries a spool of grey thread, spun from the hair of the Hollowed. When she offers you food, she ties a thread around your wrist -- invisible, weightless, unless you know to look for it. When you accept, she ties the other end to the Cord. Now you are bound. Not in a way you would notice -- no, never that. But the next time you face a choice, you will lean slightly toward the option that benefits the Chain. You will forget a name you meant to speak. You will lose a letter you meant to send. Small things. Small erasures.

Enough small things, and a kingdom falls.

\section*{How to Unbind Yourself (If You Dare)}

I write this knowing the High Mother will read it. She reads everything written in her daughters" blood, and my blood has been in this ink since I began. But perhaps someone, somewhere, will see these words before they are erased.

To unbind a Daughter's thread:

\begin{enumerate}
    \item \textbf{Find a threshold} -- a doorway, a riverbank, a crossroads at dawn. The thread is weakest where two things meet.
    \item \textbf{Light a candle of pure beeswax.} The High Mother's thread is spun from tallow; beeswax burns it away.
    \item \textbf{Speak your own name three times, aloud.} Not the name your mother gave you -- the name you have kept secret, the one you have never told anyone. That is the name the thread has stolen. You must reclaim it.
    \item \textbf{Drink a cup of salt water.} Salt is memory. The water will taste like tears. That is the taste of the cord loosening.
    \item \textbf{Do not look behind you.} The Hollowed will be there, watching. If you meet its eyes, you will forget why you began.
\end{enumerate}

If you succeed, you will feel the thread snap -- a small, sharp pain behind your left ear, like a plucked harp string. The Daughter who tied it will taste blood in her mouth for a day and a night. She will know you have broken free. She will not come after you -- the Chain does not chase. It waits. And it remembers.

\section*{The Final Confession}

I have written too much. My hand is shaking. The Hollowed in the corner of this room has not blinked in the time it has taken me to fill these pages. I know what comes next: tonight, the dream will take me, and in the morning, I will remember nothing of what I have written. The pages will be blank. My sisters will smile and say, "You look tired, Sister Alva. Rest."

But maybe -- just maybe -- someone will find these words before the erasure. Someone will read them and wonder why the milk tastes strange, why the midwife's hands are so cold, why the nurse at the hospice never blinks.

If that someone is you, do this: go home. Boil water. Salt it. Drink it. Speak your secret name into the steam. Light a beeswax candle. And never, ever accept soup from a Daughter of the Unbroken Chord.

\textit{The soup is always free. The price never is.}

\vspace{1cm}
\hfill --- \textit{Sister Alva, Weaver of the Third Spindle, Keeper of the Cord, Hollowed by choice, damned by necessity.}

\medskip
\textit{(This page was found pressed between the leaves of a herbalist's handbook in a burned cottage outside Ecktoria. The cottage had no owner on any record. The handwriting matches none of the Daughters known to the Inquisition. The Gray Wanderer refuses to discuss how they obtained it.)}

\chapter{The Quiet Work of Names -- Witches of Fate's Edge}

\epigraph{"A witch is a threshold worker---one who listens to Echo (what has been), steadies the Veil (what is), and leans into Flow (what wishes to be). Witchcraft is not brute force; it's fit: finding where the world already wants to move and giving it a gentle, precise nudge."}{\textit{--- The Gray Wanderer}}

\section{What a Witch Is (and Isn't)}

A witch is not a Runekeeper, though some keep a Codex of pressed herbs and knotted cords. Not an Invoker, though many carry a Symbol worn smooth by generations of thumb and worry. Not a Cantor, though their voice may hum old rhymes while they work. A witch practices \textit{threshold magic}: the art of finding the place where two things meet---doorway and room, waking and sleep, life and death, promise and breach---and asking that place to remember a new pattern.

There are as many kinds of witches as there are thresholds:

\begin{itemize}
    \item \textbf{Hearth Witches} who keep households safe, milk from souring, children from nightmares.
    \item \textbf{Road Witches} who know which crossroads eats liars and which ferry takes the dead.
    \item \textbf{Green Witches} who speak with nettle and briar, who taste a person's illness in the sap of a cut twig.
    \item \textbf{Bone Witches} who listen to the dead not to raise them, but to let them finish speaking.
    \item \textbf{Cord Witches} (Daughters of the Unbroken Cord) who weave obligation into flesh and memory.
\end{itemize}

Most folk magic---the kind that is indistinguishable from skilled attention---is simply witchcraft practiced well. The danger rises when a witch forces resonance instead of courting it: overbinding a vow, tightening a name, dragging one rope across many lives.

\section{First Principles}

\subsection{The Law of Named Existence}
To name is to anchor. A true name fastens an Echo; forgetting frays it. Words spoken under witness reshape how a place remembers. Every curse, blessing, or bargain is a small act of creation through recognition. Un‑Naming (suppression, forgetting) can weaken or banish.

\subsection{The Concordant Layers -- Echo / Veil / Flow}
People, places, and vows exist as chords across three layers:

\begin{itemize}
    \item \textbf{Echo:} The accumulated memory and intention that shapes reality. Stone remembers. Blood remembers. A room where a murder happened is \textit{different} afterward.
    \item \textbf{Veil:} The current, perceivable state of things; the boundary between potential and actual. This is what most people call "reality."
    \item \textbf{Flow:} The direction and potential of change, the will of elements and intent. The river wants to run downhill. A promise wants to be kept. A curse wants to execute.
\end{itemize}

Witches feel dissonance within this chord and decide whether to harmonize (benediction), redirect (hex), or lock it in place (binding). Wise law: do not sever, do not overrule, do not forget. Work with the existing harmony.

\subsection{Threshold Mechanics}
All workings start at edges---doorways, wakes, vows, last breaths, first light, crossroads, riverbanks. Magic is boundary‑law: who may pass, what must stay, and what should never meet. Thresholds are points of maximum potential for change.

\subsection{Element-Will}
Elements are wills, not mere substances. Earth (structure), Air (chance), Fate (order), Luck (opportunity), Fire (life), Water (Obishaal -- Death/Dreams), and others. A working "argues" with a chosen will; Obligation is the debt you incur when you borrow strength from a will larger than yours.

\subsection{Story Weight}
Stories have gravity. Meaning accrues; reality leans to fulfill it. When the table marks a complication (Story Beat), the world is simply cashing what the scene has earned through narrative momentum. The universe has narrative inertia.

\subsection{The Silent Ninth}
There is a missing step in the scale---acknowledged, never filled. The Ninth is absence with teeth: erasure, gaps, omissions, and the power of what is left unsaid. Tied to the Patron of the same name (and distantly to the Witness). Manipulating the Ninth is potent and dangerous.

\section{What Witches Actually Do -- Ritual Families}

Unless otherwise stated, a working requires an action (Some Time) and a roll (typically Wits + Lore or Spirit + Craft). All benefit from a \textbf{threshold}: name it, describe it, and improve Position by one step.

\subsection{Bell-Rites}
Strike, count, hold the last beat in silence---pin the Veil while you speak the Name. Tags: [WITNESS], [BIND]. Engages the layer of the Veil.
\textbf{Example:} Ring a handbell three times, then stop. In the silence, name a boundary ("no violence under this roof"). The silence witnesses the oath.

\subsection{Knot \& Cord}
Memory braided into order; untie to release a promise. Tags: [ECHO], [SEAL]. Engages Echo.
\textbf{Example:} Tie three knots in a red thread while speaking a vow. Hand the thread to a friend. When they untie it, the vow is fulfilled.

\subsection{Chalk \& Salt}
Draw a passage or a prison; shape decides which. Tags: [THRESHOLD].
\textbf{Example:} Draw a salt line across a doorway. Those you name may cross; others hesitate. The line remains until wind or water scatters it.

\subsection{Offering \& Omission}
Sometimes you add; sometimes you remove---perfect ground for Ninth‑work. Tags: [VOID].
\textbf{Example:} Pour a cup of milk onto bare earth, speaking what you wish to be forgotten. The earth drinks; the memory fades.

\subsection{Layer Focus (Choose One)}
\begin{itemize}
    \item \textbf{Echo:} +1 die to recall or restore. Backlash: recursion (loops, past actions repeat).
    \item \textbf{Veil:} +1 Position to conceal or shape. Backlash: hairline cracks (subtle reality breaks).
    \item \textbf{Flow:} +1 Effect to urge or change. Backlash: overshoot (change goes too far).
\end{itemize}

\section{Witch Paths and Talents}

A witch character is built like any other (attributes, skills), but replaces standard magic with \textbf{Threshold Talents}. These talents are available to any character who learns witchcraft (often via a mentor, a grimoire, or a Patron's dream). The three branches---Hearth, Passage, Memory---can be advanced independently.

\subsection{Hearth Branch -- Protection \& Sustaining}

\textbf{Warm Hand} (2 XP). When you tend someone in need (food, comfort, shelter), gain +1 die to your next Threshold working in that space.

\textbf{Cup and Salt} (2 XP) -- Prereq: Warm Hand. Once per session, declare an area under Guest‑Rite: violence suffers −1 Effect; negotiation gains +1 Position.

\textbf{Hearth's Memory} (3 XP) -- Prereq: Cup and Salt. A place you have protected remembers you. Once when returning, choose: Position +1 for a scene, a clue revealed, or hidden harm heals 1.

\textbf{Keeper of Thresholds} (4 XP) -- Prereq: Hearth's Memory. Name a dwelling or crossing as your ward. Once per scene you can: reduce a community clock by 1, or upgrade Guest‑Rite to Sanctuary (no initiations of violence).

\subsection{Passage Branch -- Movement \& Mediation}

\textbf{Step Between} (2 XP). Once per scene, ignore a Position penalty due to contested faith, order, or personal misalignment at a threshold.

\textbf{Bridge‑Maker} (3 XP) -- Prereq: Step Between. When mediating between sides at a crossing or doorway, you and one ally gain +1 die to social rolls for the scene.

\textbf{Traveler's Favor} (3 XP) -- Prereq: Bridge‑Maker. Once per journey (or per session), treat difficult border or bureaucracy as one step easier (DV −1 or Position +1).

\textbf{Keeper of Crossroads} (4 XP) -- Prereq: Traveler's Favor. Once per session, declare a Crossroads Binding: until scene end, no party may leave without naming their choice -- and consequences adjust accordingly.

\subsection{Memory Branch -- Story, Names, \& Reciprocity}

\textbf{Soft Witness} (2 XP). You count as witness for your own work. When another watches in respect, you and they gain +1 die for the scene.

\textbf{Name‑Keeper} (3 XP) -- Prereq: Soft Witness. When you learn a true name, bind it to a memento. Once per session, invoke it to improve Position +1 with that person or gain a clue about them.

\textbf{Debt‑Binder} (3 XP) -- Prereq: Name‑Keeper. Once per scene, you may mark or clear one segment of a Promise Clock when you resolve or escalate a conflict through story, hospitality, or grief.

\textbf{Weaver of Remembering} (4 XP) -- Prereq: Debt‑Binder. Once per session, you may: rewrite the social consequence of a working, convert backlash into obligation, or create a Legacy Blessing that gives +1 die in a future scene tied to this threshold.

\subsection{Capstones (Choose one per tree)}

\textbf{House‑Mother / House‑Father} (6 XP). Consecrate a dwelling or camp: Threshold workings gain +1 Effect within it, backlash softens into obligation instead of harm, and the house becomes a named NPC with disposition.

\textbf{Road‑Saint} (6 XP). Treat boundaries as your domain. Once per scene, declare: safe passage, stalled pursuit, or a turning of fate (reroll one ally or enemy pool).

\textbf{Saga‑Bearer} (6 XP). Your words reshape identity. When you speak someone's forgotten truth or name: they become an NPC ally for the scene, or if enemy, they hesitate or falter (downgrade Harm or Condition by 1).

\section{Patrons of the Witches}

Most witches do not serve a single Patron, but those who do (sometimes called warlocks) gain focused gifts. These Patrons are particularly sympathetic to threshold magic.

\subsection{Inaea -- Angel of the Spider (Mercy, Continuity)}
\textbf{Gift:} Mercy Lines. Once per scene, you may declare a harm stopped at a door‑thread. An ally in Near ignores the next Harm 1.
\textbf{Backlash:} Mercy without memory hollows. Apologies left unsaid become cold rooms -- mark a Promise Clock [4] that ticks when you fail to witness a pledge.

\subsection{Isoka -- Angel of Serpents (Change, Shedding)}
\textbf{Gift:} Loosening Skins. Once per scene, unhook a role or label from yourself or an ally: remove one Condition or downgrade Harm 1 to Fatigue.
\textbf{Backlash:} Purges; identity slips; loyalties molt too. Mark 1 SB (Hearts) when you shed something that should have been kept.

\subsection{Ikasha -- She Who Sleeps (Shadow, Possibility)}
\textbf{Gift:} Shadow Courtesy. Once per scene, declare a night crossing held safe by role and bell. You and allies gain +1 die to Stealth for the scene.
\textbf{Backlash:} Too many possible selves peer in. At scene end, test Resolve (DV 3) or suffer Disorient (Fatigue 1).

\subsection{Morag the Hag (Bargains, Hidden Costs)}
\textbf{Gift:} The Crooked Thread. When you bind a minor promise in red thread, the target suffers −1 die to break it. The thread lasts one scene.
\textbf{Backlash:} Compulsive Bargaining. You must name a price whenever value changes hands. If you refuse, mark 1 SB (Diamonds).

\subsection{Lunara -- The Silver Quiet (Moon, Mystery)}
\textbf{Gift:} Moonlit Mirror. Once per session, gaze into any reflective surface under moonlight. Ask one question about a hidden truth; the GM answers truthfully (perhaps cryptically).
\textbf{Backlash:} Shadows cling. Your eyes glimmer in dim light, unnerving others -- −1 die to social rolls with the suspicious.

\subsection{Nimorith -- The Grey Benefactor (Absence, Threshold)}
\textbf{Gift:} Shadow‑Step. Once per scene, move from one witnessed threshold to another within Near. A door you can see, a window, a bell's shadow.
\textbf{Backlash:} Almost‑Presence. In crowds you slip notice; begin Position −1 when trying to stand out.

\subsection{The High Mother (Unbroken Chord)}
\textbf{Gift:} Cord of Flesh. You may store up to Body rating points of Fatigue as Chained Knots, appearing uninjured. At sunset, you must sacrifice or take 2 Harm.
\textbf{Backlash:} The Hollowed watches. Your Hollowed familiar (a heartless child or animal) never blinks. Others notice something wrong −1 die to social rolls with the pious.

\section{Curses as Architecture (Keeper Tools)}

A curse is an edit to the memory rulebook---a habit inserted into how the world recalls a person or object. Write a curse like a rule:

\emph{"Until a bell is rung in your true name, lamps forget to light for you."} (Veil‑edit, Named anchor)

\subsection{Curses have a Clock}
\textbf{Recursion [6]} -- ticks when the victim reinforces the story (fumbling, hiding the problem, refusing witness).

\subsection{Provide a counter}
Add the missing layer: Name + Witness + Threshold. Example: ring the bell at a doorway with kin present.

\subsection{Story Beat Menu for Curses}
\begin{itemize}
    \item 1 SB: minor echo (old order resurfaces)
    \item 2 SB: warped threshold (wrong door opens)
    \item 3 SB: layer slip (Veil cracks, Flow surges)
    \item 4 SB: Ninth propagation (a second gap appears)
\end{itemize}

\section{Witch-Hunters \& Containment Orders}

Every culture that lives beside Patrons and old ways develops a method to contain what runs wild. Some orders prosecute crimes; others prevent collapses of context. Witch‑hunters are often as complex as the witches they pursue.

\subsection{Chain‑Lanterns of Thepyrgos (Human Witch Inquisitors)}
Lanterns hooded with iron mesh. They require a witness before any accusation. Mark a Lantern Line (women's rites cannot cross without test), issue a Father's Claim to bind property, serve a Stop‑Rite halting a working until patriarchal review.

\subsection{Spirit‑Shield Correctors (Aeler)}
Counting breaths at thresholds; bells on lintels, always rung by fathers. They seal rogue workings into Repair Ledgers -- convert unlicensed magic into domestic labor.

\subsection{Break‑Reins Riders (Ykrul)}
Unknotted bridles on cairns. They cut coercive bindings, escort witches under guest‑right. No rope becomes a leash.

\subsection{Silk Vigil (Lethai‑ar, Inaea/Isoka)}
Three‑tone bells; hair‑fine lines across thresholds. They enforce Mask Right (declare roles, ring once for guesthood). Wedding Line sanctuary; Context Reset (restore roles, end duels of pride).

\subsection{Running a Witch Hunt -- Quick Clocks}
\begin{itemize}
    \item \textbf{Purity Panic [4]} -- public fear escalates.
    \item \textbf{Due Process [4]} -- legal constraints bind hunters.
    \item \textbf{Hospitality vs. Sanction} -- if the community feeds the accused, hunters lose a die.
\end{itemize}
\textbf{Resolution:} On a successful Containment, the rogue magic is boxed or blessed. On failure, it spreads to a new threshold.

\section{Sample Rites for Player Witches}

\subsection*{Bless Hearth (Low, DV 3)}
\textit{"Salt in the corners, milk on the stone, fire in the hearth, you are home."}
Components: Salt, a cup of milk, a lit candle.
Effect: For one night, those who sleep within the blessed space recover Fatigue twice as fast, and the first intruder entering with ill intent becomes immediately known to the caster (a sudden chill, a candle guttering). The blessing ends if the milk is spilled or the fire dies.

\subsection*{Name the Unspoken Grief (Standard, DV 4)}
\textit{"What was not said, speak now through this knot."}
Components: A red thread, a personal token of the grieving.
Effect: Draw the thread through the token while whispering the grief you have witnessed. The thread knots itself once for each untold sorrow. Thereafter, the grieving may speak of that pain without choking, and you gain +1 die to aid them. If you break the thread before the grief is fully borne, the sorrow returns thrice worse.

\subsection*{Bind the Wanderer (Standard, DV 4)}
\textit{"By this door, by this stone, by this name you have known -- stay."}
Components: A door threshold, a name scratched on a clay shard.
Effect: One supernatural entity (ghost, Trow, spirit) cannot cross the named threshold until dawn, unless you willingly break the shard. The binding is peaceful -- the entity may watch, but not enter. Breaking the shard early marks 1 SB (Clubs) as the threshold cracks.

\subsection*{The Missing Line (High, DV 5)}
\textit{"Let this line be unlived. Let this word be unwitnessed."}
Components: A document or contract, a grey feather that casts no shadow.
Effect: Cause one clause, name, or line to vanish from a witnessed text. The erasure is seamless and socially "always so." The missing information persists in Echo but not in Veil -- it may be rediscovered by deep investigation. Obligation: 2 segments.

\section{Cottagecraft -- Everyday Magic}

Witchcraft shines in the small moments. These minor workings require no roll, only appropriate components and a few minutes of quiet.

\begin{itemize}
    \item \textbf{Keep Milk from Souring:} Stir thrice with a hawthorn twig. The milk stays sweet for a week.
    \item \textbf{Find a Lost Object:} Tie a knot in a handkerchief while speaking the object's name. Walk the room; the knot will tug when you are near.
    \item \textbf{Silence a Squeaking Door:} Anoint the hinge with a drop of olive oil, whispered, "You have spoken enough."
    \item \textbf{Remember a Dream:} Place a piece of amethyst or a smooth river stone under your pillow. Recite your name backward before sleep.
    \item \textbf{Ward a Threshold against Nightmares:} Slip three bay leaves and a pinch of salt beneath the doormat. Renew at each new moon.
\end{itemize}

\section{The Gray Wanderer's Last Advice on Witches}

\epigraph{"The best witches I have known were not powerful. They were patient. They knew which knot to tie and which to leave undone, which name to speak and which to forget. That is the quiet work of names -- and it is enough to change the world, one threshold at a time."}{\textit{--- The Gray Wanderer}}

\chapter{The Mind Ascendant -- Psionics and the Gray Wanderer}

\epigraph{"I have seen a man stop a drawn sword with his forehead. I have watched a woman read a sealed letter without breaking the wax. I have felt a child's hand on my shoulder when no child was there. And I still cannot tell you how any of it works. Magic? I understand magic. Debts, Patrons, the weight of a broken vow. But this? This is something else entirely."}{\textit{--- The Gray Wanderer}}

\section{The Wanderer's Confession -- A Conversation at the Dusty Page}

I met the monk at a crossroads inn, the sort where the floorboards are older than the kingdom and the ale tastes faintly of regret. He wore the grey robes of an ascetic from the Vhasian badlands -- not a warrior-monk with shaved head and calloused fists, but a \textit{meditant}, a man who had spent forty years learning to sit still and listen to the shape of his own thoughts.

"You"re a psion," I said. It was not a question.

He smiled. "I am a \textit{tongue of the silent mind}. But yes. You name it psionics. The Aeler call it "deep accounting." The Ykrul call it "weather-reading in the skull." The Lethai call it "hearing without a bell.""

"And what do you call it?"

"Practice."

I ordered another drink. "I"ve spent my life bargaining with Patrons, carving Rites into my own skin, singing songs that should not exist. I understand power. But psionics? It has no Patron. No Symbol. No Codex. No Thiasos. No Obligation. How can something that powerful have \textit{no cost}?"

The monk laughed -- a dry, sandpaper sound. "No cost? Wanderer, you mistake the nature of the debt. A Runekeeper owes a Patron. A Cantor owes their voice. A psion owes \textit{themselves}. Every thought you bend is a thought you cannot unthink. Every mind you touch leaves a fingerprint on your own. Every future you glimpse steals a moment of your present. The cost is not external. It is \textit{attrition}. The slow wearing away of the self."

He tapped his temple. "This is the only ledger that matters. And it always comes due."

\section{What Psionics Is (and What It Is Not)}

Psionics is the disciplined mastery of the mind's hidden potentials. Unlike magic -- which borrows power from Patrons, elements, or the Weave -- psionics is \textit{autonomous}. The power comes from within, shaped by will and honed by practice. This makes psions both powerful and vulnerable: their greatest strength is also their most dangerous weakness.

\textbf{Psionics is not:}
\begin{itemize}
    \item \textbf{Magic.} It does not invoke Patrons, elements, or the Ninth. Anti-magic wards do not stop it -- but \textit{psychic dampening fields} do.
    \item \textbf{Innate.} Psionics requires training, discipline, and often physical or meditative practice. A natural talent is like a sharp knife -- useful, but dangerous to the untrained hand.
    \item \textbf{Safe.} Mental Strain is real. Mental scars are permanent. The mind that bends reality also bears the cracks.
\end{itemize}

\textbf{Psionics is:}
\begin{itemize}
    \item \textbf{Internal.} A psion's power is measured by Mental Strain, a resource that recovers with rest and discipline.
    \item \textbf{Subtle.} Psionic effects often have no visible component -- a thought, a glance, a held breath. This makes psions excellent investigators, infiltrators, and spies.
    \item \textbf{Consequential.} Every psionic action generates Story Beats. The universe notices when a mind touches what it should not.
\end{itemize}

\section{Core Mechanics -- The Mind's Ledger}

\subsection{Mental Strain Track}
Mental Strain represents your capacity for sustained psychic exertion.

\[
\text{Mental Strain Capacity} = \text{Spirit}\ (\text{minimum } 2)
\]

Mark Mental Strain when you use psionic Arts, resist psychic intrusions, or push your mind beyond its limits. When you cannot pay a Mental Strain cost, choose:
\begin{itemize}
    \item Convert the unpaid point to \(+2\) regular Fatigue, or
    \item Convert the unpaid point to \(+1\) Harm (Stress).
\end{itemize}

\textbf{Overflow:} If you ever exceed your Mental Strain capacity, immediately suffer \(+1\) Harm (Stress) and set Mental Strain to your maximum (no further over-marking this scene).

\subsection{Psionics Skill}
The Psionics skill (0--5) gates access to psionic Arts and sets control and efficiency. You cannot use psionic Arts without investing in this skill. Each Art has a \textbf{Key Attribute} (see table). Using a non-key attribute imposes \(+1\) DV or \(-1\) die (GM's choice).

\subsection{Resolution}
Roll \textbf{Attribute + Psionics} vs. DV. The GM sets DV based on scope, risk, and opposition (typical 2--5+). Each 1 on the roll generates a Story Beat as normal.

\section{Psionic Arts}

\subsection{Telekinesis}
\textbf{Key Attribute:} Spirit. \textbf{Alternate:} Wits (+1 DV for fine manipulation). \textbf{Action:} 1 action; sustain by fiction. \textbf{Range:} Near (+1 DV for Far). \textbf{Cost:} 1 Mental Strain to activate; sustained lifts/holds cost 1 Mental Strain per scene.

\textbf{SB:} Scales with mass/force:
\begin{itemize}
    \item \textbf{Light (DV 2--3):} cups, tools, doors. No SB.
    \item \textbf{Medium (DV 4):} furniture, a person's mass, riding debris. +1 SB on Partial/Miss.
    \item \textbf{Heavy (DV 5+):} carts, statues, boulders, multi-target shoves. +2 SB on Partial/Miss; +1 SB on Success.
\end{itemize}

\textbf{Uses:} Hurl objects as weapons, shape cover, remote tasks. Cannot directly crush living tissue -- injury is environmental (falling, debris).

\subsection{Telepathy}
\textbf{Key Attribute:} Wits. \textbf{Alternate:} Presence (+1 DV, emotive broadcast). \textbf{Action:} 1 action. \textbf{Range:} Near, line of sight (+1 DV for Far; +1 DV without LoS if strong link). \textbf{Targets:} One (+1 DV per extra target this scene).

\textbf{Depth \& SB -- choose before rolling:}
\begin{itemize}
    \item \textbf{Surface Thoughts (DV 2):} current focus. No SB first time per target/scene; repeats +1 DV, on Miss +1 SB.
    \item \textbf{Emotional State (DV 3):} tones/urges. +1 SB on Partial/Miss.
    \item \textbf{Deep Memories (DV 5):} episodic recall. +1 SB on Success; +2 SB on Partial/Miss.
\end{itemize}

\textbf{Outcomes:} Success grants information and one set‑up use (angle, position, pointed question). Partial gives blurred fragments. Miss gives no read; target feels the brush.

\textbf{Defense:} Unwilling targets roll Spirit + Resolve; if they meet or beat your roll, your read fails and you mark +1 SB (+2 for Deep). Margin 2+ reveals your intent.

\subsection{Clairvoyance}
\textbf{Key Attribute:} Wits. \textbf{Alternate:} Spirit (+1 DV, trance method). \textbf{Action:} 1 action; ritual scry may take a scene. \textbf{Range:} Scene/locale by fiction; +1 DV per range band beyond Near. \textbf{Cost:} DV 2--3 = 1 Mental Strain; DV 4 = 2; DV 5+ = 3.

\textbf{SB:} Near none; Distant +1 SB on Partial/Miss; Complex Scry +2 SB on Partial/Miss; +1 SB on Success. Covert spying on private moments generates +1 SB.

\subsection{Biofeedback}
\textbf{Key Attribute:} Spirit. \textbf{Alternate:} Wits (+1 DV, clinical technique). \textbf{Action:} 1 action (major work takes a focused scene). \textbf{Intensity:}
\begin{itemize}
    \item \textbf{Minor (DV 2):} stabilize, close wounds, purge shock. No SB. Heal up to Psionics levels per day (cap 1 Harm per recipient/day).
    \item \textbf{Moderate (DV 3):} +1 SB on Partial/Miss; heal up to Psionics (max 2/day).
    \item \textbf{Major (DV 4+):} +2 SB on Partial/Miss; +1 SB on Success; heal up to Psionics (max 1/day); mark +1 Fatigue.
\end{itemize}

\textbf{Perks:} Convert 1 Harm ↔ 1 Mental Strain (up to Psionics, max 3/day). Brief physical enhancement (+1 die).

\subsection{Astral Projection}
\textbf{Key Attribute:} Spirit. \textbf{Alternate:} Wits (+1 DV, skim tether). \textbf{Action:} 1 action to slip; travel by fiction. \textbf{Cost:} 1 Mental Strain per scene. \textbf{SB:} +1 per band beyond Near; +1 per additional scene sustained.

\textbf{Risks:} Body is Helpless. Astral Harm echoes to flesh. Wards add +2 DV. Extended projection (3+ scenes) courts separation; astral death = physical death.

\subsection{Psychic Assault}
\textbf{Key Attribute:} Spirit. \textbf{Alternate:} Wits (+1 DV, needle‑strike). \textbf{Action:} 1 action. \textbf{Cost:} DV 2--3 = 1 Mental Strain; DV 4 = 2; DV 5+ = 3. \textbf{SB:} always +1, plus DV 4--5 additional.

\textbf{Effect:} Mental Harm, bypasses physical armor. Can target specific faculties (sight, balance, focus) by fiction. High‑Spirit foes may add +1 DV.

\subsection{Mind Shield}
\textbf{Key Attribute:} Wits. \textbf{Alternate:} Spirit (+1 DV, soak‑and‑ground). \textbf{Action:} 1 action. \textbf{SB:} Passive none; per effect blocked +1 SB; Area +2 SB and mark +2 Fatigue.

\textbf{Effect:} Resist intrusion/telepathy; shelter allies; create dampened zones. Protecting others costs 2× Mental Strain. Area protection covers all within Near (friend and foe).

\subsection{Empathic Manipulation}
\textbf{Key Attribute:} Presence. \textbf{Alternate:} Wits (+1 DV, cold‑read leverage). \textbf{Action:} 1 action. \textbf{Intensity:}
\begin{itemize}
    \item \textbf{Subtle (DV 2--3):} nudge mood, steady hands. No SB.
    \item \textbf{Moderate (DV 4):} calm rage, spark courage. +1 SB on Partial/Miss.
    \item \textbf{Strong (DV 5+):} sway a crowd's tenor. +2 SB on Partial/Miss; +1 SB on Success.
\end{itemize}

\textbf{Risks:} You feel what you project. Same target not more than once per scene. Opposing core values: +2 DV. Affecting a crowd always generates +1 SB.

\subsection{Precognition}
\textbf{Key Attribute:} Spirit. \textbf{Alternate:} Wits (+1 DV, tactical glimpse). \textbf{Action:} 1 action; deeper visions may take a scene. \textbf{Cost:} DV 2--3 = 1 Mental Strain; DV 4 = 2; DV 5+ = 3. \textbf{SB (inherent):} Minor +1; Major (DV 4) +2; Detailed (DV 5+) +3.

\textbf{Effect:} Advantage on future rolls, avoid named dangers, compare probable branches. \textbf{Paradox -- Future Fixation:} Until downtime (or a defiant choice), lose 1 Boon per session; actions that resist the shown path may be Desperate by GM call.

\subsection{Base Costs Summary}
\begin{itemize}
    \item Activation: 1 Mental Strain
    \item Maintenance: 1 Mental Strain per scene / Significant Time
    \item Maintenance Reduction: Wits + Psionics (DV 3) reduces cost by 1; failure generates 1 SB.
\end{itemize}

\section{Psionic Talents}

\subsection{Minor Talents (2--4 XP)}
\begin{itemize}
    \item \textbf{Mental Fortress (3 XP):} +1 die to resist psychic effects; once per scene, convert incoming psychic Harm to Fatigue.
    \item \textbf{Thought Thief (4 XP):} +1 die to read surface thoughts; once per scene, detect lies automatically.
    \item \textbf{Psychic Reservoir (4 XP):} Increase Mental Strain track by +2; recover 1 additional Mental Strain during rest.
\end{itemize}

\subsection{Major Talents (6--8 XP)}
\begin{itemize}
    \item \textbf{Telekinetic Mastery (7 XP):} +1 die to Telekinesis attacks; lift objects one category heavier without increasing DV.
    \item \textbf{Mind Walker (8 XP):} Astral Projection requires no concentration; move at normal speed in astral form; reduce projection cost to \(1/2\) Mental Strain (round up).
    \item \textbf{Psychic Vampire (8 XP):} When dealing psychic Harm, heal 1 Mental Strain; once per scene, drain 1 Mental Strain from target (requires successful Psychic Assault).
\end{itemize}

\subsection{Prestige Talents (12+ XP)}
\begin{itemize}
    \item \textbf{Master of Minds (12 XP):} +2 dice to all telepathic effects; once per scene, impose emotional state on a group (DV 4, mark 2 Fatigue).
    \item \textbf{Mind Over Body (12 XP):} Mental Strain track becomes Spirit + Wits (min 6, cap 12). When you would suffer Mental Strain overflow, convert each overflow point into +1 Fatigue instead.
    \item \textbf{Reality Bender (15 XP):} Once per session, reshape local reality within Near range. Minor alteration (DV 3), moderate (DV 4), major (DV 5). Mark 1 Harm and 1 Fatigue; generate SB based on severity.
\end{itemize}

\section{The Great Silence -- A History the Wanderer Missed}

The monk leaned forward. "You want to understand psionics? First you must understand what was \textit{lost}."

He spoke of the Great Silence -- a period two to three centuries long, thousands of years ago, when psionic abilities across all peoples were suppressed, compressed, or transformed. What had once been natural gifts became rare talents requiring intense training and carrying significant risk.

"The Aeler sealed their psychic academies," he said. "The Ykrul buried their "mind roads." The Lethai forgot how to hear without bells. And when the Silence lifted -- in cycles, each Awakening bringing back different abilities -- the powers had changed. They were harder. More costly. More dangerous."

\textbf{Resonance Scars:} Certain locations still bear the marks of the Great Silence:
\begin{itemize}
    \item \textbf{The Whispering Vaults (Aeler Holds):} Amplify Mental Strain by 2 segments but grant +1 Effect to telepathy.
    \item \textbf{The Empty Circle (Ykrul Steppe):} No psionic abilities function at all.
    \item \textbf{The Singing Stones (Lethai Valewood):} Allow long‑distance telepathic communication when properly "tuned."
\end{itemize}

\textbf{Awakening Cycles:} Since the Silence, psionic abilities have returned in waves:
\begin{itemize}
    \item \textbf{First Awakening (150 years ago -- present):} Telepathy and telekinesis return.
    \item \textbf{Second Awakening (400--325 years ago):} Precognition and biofeedback appear.
    \item \textbf{Third Awakening (650--575 years ago):} Astral projection and empathic manipulation.
\end{itemize}

"We are in the First Awakening now," the monk said. "What abilities will come next -- and what they will cost -- no one knows."

\section{The Thought Trade -- Psionic Economics}

"You think psionics is only for battle and spying?" The monk shook his head. "In the free ports, psionic labor is a commodity like grain or timber."

\subsection{Memory Merchants}
Telepaths who store and retrieve memories for merchants, scholars, and nobles. A secure memory storage hour costs a skilled craftsman's daily wage. Dangerous memories cost ten times that.

\subsection{Calculation Circles}
Psions with enhanced mathematical abilities act as living computers for complex engineering and navigation. They perform in hours what takes traditional mathematicians months -- but mental strain requires significant recovery.

\subsection{Experience Exchange}
In underground markets, psions trade raw experiences -- memories of adventure, intense emotions, unique sensations -- as currency. A dragon-flight memory might be worth fifty silver. The experience of true love? Several gold.

\subsection{Skill Debt}
Psions enter arrangements where they provide mental services over time in exchange for immediate payment. This creates a complex system of mental credit and obligation, with contracts enforced by guilds or, in less regulated areas, reputation and threat.

\section{Psionic Investigation -- The Detective's Third Eye}

"A murder happens," the monk said. "No witnesses, no evidence. A traditional investigator has nothing. A psion has a \textit{residual impression} -- the echo of the last moments, pressed into stone and shadow like a fingerprint."

\textbf{Residual Impression Reading:} Wits + Psionics vs. DV 4--6 (time elapsed, emotional intensity, interference). Success grants vivid impressions of final moments. Partial gives fragmented, potentially misleading fragments. Miss creates Psychic Contamination -- the psion absorbs traumatic elements that affect future investigations.

\textbf{Legal Complications:} Most jurisdictions struggle with mental evidence. Aeler courts require dual testimony (psionic + traditional). Lethai moots accept mental evidence only with proper context witnesses. Ykrul bowl proceedings accept it if obtained without deception. Free havens are more permissive but require the psion to testify about methods and error potential.

\section{The Monk's Final Words}

The inn was quiet. The fire had burned low. The monk stood, draped his grey robes over his shoulders, and placed a smooth river stone on the table between us.

"Take this," he said. "It is a resonator. When you touch it, you will feel the shape of your own thoughts more clearly. Use it to meditate. Use it to listen. But do not use it to reach into minds that have not invited you. That is the first law of the silent mind: \textit{consent is not optional.}"

I picked up the stone. It was warm. I could feel my own pulse in it, and something else -- a faint hum, like a tuning fork struck in another room.

"What is your order called?" I asked.

He smiled. "We have no name. Names are thresholds, and we prefer to remain... uninvited."

He walked out into the rain. His footsteps faded. The stone cooled.

And I sat there, for the first time in decades, wondering if I had finally met something I could not explain away with Patrons and Rites and the weight of broken vows.

\epigraph{"The mind is a door. Psionics is the key. But not every door should be opened -- and not everyone who knocks should be answered."}{\textit{--- The Gray Wanderer, from a letter never sent}}

\appendix
\AtBeginEnvironment{longtable}{\small}

\chapter{Comparing the Paths of Power -- A Wanderer's Cheat Sheet}

\epigraph{"I have walked every road in this book. Some are paved, some are thorns, some are just a goat track and a prayer. Here is what I wish someone had told me before I started."}{\textit{--- The Gray Wanderer}}

\section*{How to Read This Table}

The paths of power in Fate's Edge are not better or worse -- they are \textit{different}. Each serves a different kind of player and a different kind of story. This table compares them across several axes:

\begin{itemize}
    \item \textbf{Source:} Where the power comes from (Patron, self, spirits, etc.).
    \item \textbf{Cost:} What you pay each time you use it (Obligation, Mental Strain, Leash ticks, etc.).
    \item \textbf{Complexity:} Low / Medium / High -- how many moving parts you need to track.
    \item \textbf{Narrative Weight:} How much story emerges naturally from using the power.
    \item \textbf{Flexibility:} How many different problems the path can solve.
    \item \textbf{Risk of Burnout:} How likely you are to cripple yourself.
    \item \textbf{Best For:} Player archetypes or campaign styles.
\end{itemize}

\medskip
\noindent \textit{Note:} These are subjective -- your table may find different strengths and weaknesses. The Gray Wanderer reserves the right to be wrong.

\vspace{0.5cm}

% Replace the "Magic Systems Compared" table with this:
\small
\begin{longtable}{p{2cm} p{2cm} p{1.5cm} p{1.5cm} p{1.5cm} p{1.5cm} p{2cm}}
\caption{Magic Systems Compared}\\
\toprule
\textbf{Path} & \textbf{Source} & \textbf{Cost} & \textbf{Complexity} & \textbf{Narrative} & \textbf{Flex} & \textbf{Best For} \\
\midrule
\endfirsthead
\multicolumn{7}{c}{{\bfseries \tablename\ \thetable{} -- continued}} \\
\toprule
\textbf{Path} & \textbf{Source} & \textbf{Cost} & \textbf{Complexity} & \textbf{Narrative} & \textbf{Flex} & \textbf{Best For} \\
\midrule
\endhead
\bottomrule
\endfoot
Runekeeper & Patron, Codex & Obligation & Med & High & Low & Structure lovers \\
Invoker & Symbols & Obligation + Cond. & Med--High & High & Med & Gamblers \\
Cantor & Voice, Audience & Corruption & Low--Med & High & Med & Performers \\
Summoner & Spirits & Leash & Med & High & Low--Med & Tactical players \\
Free Caster & Raw Weave & Backlash & Low & Med & Very High & Improvisers \\
Witch & Thresholds & Promise Clocks & Low--Med & Very High & Med & Storytellers \\
Psion & Internal & Mental Strain & Med--High & High & High & Resource managers \\
\bottomrule
\end{longtable}
\normalsize

\section*{Power Source Deep Dive}

% Replace the "Magic Systems Compared" table with this:
\small
\begin{longtable}{p{2cm} p{2cm} p{1.5cm} p{1.5cm} p{1.5cm} p{1.5cm} p{2cm}}
\caption{Magic Systems Compared}\\
\toprule
\textbf{Path} & \textbf{Source} & \textbf{Cost} & \textbf{Complexity} & \textbf{Narrative} & \textbf{Flex} & \textbf{Best For} \\
\midrule
\endfirsthead
\multicolumn{7}{c}{{\bfseries \tablename\ \thetable{} -- continued}} \\
\toprule
\textbf{Path} & \textbf{Source} & \textbf{Cost} & \textbf{Complexity} & \textbf{Narrative} & \textbf{Flex} & \textbf{Best For} \\
\midrule
\endhead
\bottomrule
\endfoot
Runekeeper & Patron, Codex & Obligation & Med & High & Low & Structure lovers \\
Invoker & Symbols & Obligation + Cond. & Med--High & High & Med & Gamblers \\
Cantor & Voice, Audience & Corruption & Low--Med & High & Med & Performers \\
Summoner & Spirits & Leash & Med & High & Low--Med & Tactical players \\
Free Caster & Raw Weave & Backlash & Low & Med & Very High & Improvisers \\
Witch & Thresholds & Promise Clocks & Low--Med & Very High & Med & Storytellers \\
Psion & Internal & Mental Strain & Med--High & High & High & Resource managers \\
\bottomrule
\end{longtable}
\normalsize

\section*{Risk and Consequences}

\noindent \textit{How each path can ruin your day (and make the story better).}

% Replace the "What Fuels Your Magic" table with this:
\small
\begin{longtable}{p{2.5cm} p{5cm} p{4cm}}
\caption{What Fuels Your Magic}\\
\toprule
\textbf{Path} & \textbf{You Borrow From} & \textbf{You Owe} \\
\midrule
Runekeeper & Single Patron's Rites & Obligation → Patron Intrusion \\
Invoker & Multiple Patrons (Symbols) & Obligation + Symbol condition \\
Cantor & Patron + audience & Corruption Clock → Patron burden \\
Summoner & The spirit itself & Leash → spirit acts (then departs) \\
Free Caster & Raw Weave & Backlash (SB) → complications \\
Witch & Thresholds (Echo/Veil/Flow) & Promise Clocks → price comes due \\
Psion & Your own mind & Mental Strain → overflow = Harm \\
\bottomrule
\end{longtable}
\normalsize

% Replace the "What Goes Wrong" table with this:
\small
\begin{longtable}{p{2.5cm} p{4cm} p{4.5cm}}
\caption{What Goes Wrong (Interesting Consequences)}\\
\toprule
\textbf{Path} & \textbf{Typical Failure} & \textbf{Interesting Consequence} \\
\midrule
Runekeeper & Obligation exceeds capacity & Patron Intrusion mid-scene \\
Invoker & Symbol compromised/shattered & Crack the Seal breaks Symbol \\
Cantor & Corruption Clock fills & Permanent Patron burden \\
Summoner & Leash fills & Spirit acts to its nature (chaos) \\
Free Caster & Backlash (SB spent) & Weave pushes back (wrong effects) \\
Witch & Promise Clock fills & Neighbor demands payment \\
Psion & Mental Strain overflow → Harm & Mental scars, contamination \\
\bottomrule
\end{longtable}
\normalsize

\section*{Narrative vs. Mechanical -- Where Each Path Shines}

% Replace the "Lean Into the Fiction" table with this (split into two tables):
\small
\begin{longtable}{p{2.5cm} p{5cm}}
\caption{Lean Into the Fiction -- Mechanical Strengths}\\
\toprule
\textbf{Path} & \textbf{Mechanically Strong At} \\
\midrule
Runekeeper & Reliable, repeatable effects. Predictable costs. \\
Invoker & Versatility (multiple Patrons). Emergency power. \\
Cantor & Crowd control. Ally buffs. Emotional manipulation. \\
Summoner & Action economy (spirit acts separately). Information. \\
Free Caster & Absolute flexibility. Problem-solving creativity. \\
Witch & Social magic. Hospitality laws. Curses as architecture. \\
Psion & Subtlety. Mental investigation. Resource management. \\
\bottomrule
\end{longtable}

\vspace{0.5cm}

\begin{longtable}{p{2.5cm} p{5cm}}
\caption{Lean Into the Fiction -- Narrative Strengths}\\
\toprule
\textbf{Path} & \textbf{Narratively Strong At} \\
\midrule
Runekeeper & Patron relationships. Weight of covenant. \\
Invoker & Desperation. Stolen fire. Gambling with power. \\
Cantor & Performance. Cult leadership. Cost of fame. \\
Summoner & Cross-cultural worldbuilding. Bargaining. \\
Free Caster & Raw, untamed magic. Backlash as story engine. \\
Witch & Folk horror. Community ties. Quiet work of names. \\
Psion & Internal struggle. Cost of enlightenment. Being hunted. \\
\bottomrule
\end{longtable}
\normalsize

\section*{The Wanderer's Recommendations}

\noindent \textit{If you are a player who...}

\begin{itemize}
    \item ...loves structure, ledgers, and owing a powerful boss: play a \textbf{Runekeeper}.
    \item ...wants to live dangerously and flip between masters: play an \textbf{Invoker}.
    \item ...wants to be the face of the party and sing for your supper: play a \textbf{Cantor}.
    \item ...likes having a friend (who might eat you) and exploring strange cultures: play a \textbf{Summoner}.
    \item ...hates spell lists and wants to invent effects on the fly: play a \textbf{Free Caster} (with the Spellcraft talent).
    \item ...wants magic that feels like gardening, negotiation, and quiet rebellion: play a \textbf{Witch}.
    \item ...wants internal resource management, no Patron, and mind‑bending themes: play a \textbf{Psion}.
\end{itemize}

\noindent \textit{If you are a GM looking for...}

\begin{itemize}
    \item ...a low‑complexity magic user for a new player: \textbf{Cantor} or \textbf{Witch} (basic rites).
    \item ...a high‑complexity system for a veteran: \textbf{Invoker} or \textbf{Psion}.
    \item ...magic that generates story hooks automatically: \textbf{Runekeeper} (Patron Intrusions) or \textbf{Witch} (Promise Clocks).
    \item ...magic that feels alien and dangerous: \textbf{Free Caster} (Backlash) or \textbf{Summoner} (Leash).
    \item ...magic that integrates with social and investigative scenes: \textbf{Witch} or \textbf{Psion}.
    \item ...magic that is flashy and public: \textbf{Cantor}.
    \item ...magic that is subtle and secret: \textbf{Psion} or \textbf{Witch} (Ninth‑work).
\end{itemize}

\section*{A Final Note on Mixing Paths}

Nothing prevents a character from walking multiple roads -- a Runekeeper who also dabbles in free casting, a Witch who summons spirits, a Psion who sings Cantor's hymns. However, each path tracks its own costs separately (Obligation, Corruption, Mental Strain, Promise Clocks, Leash, Backlash). Mixing paths multiplies bookkeeping. The Gray Wanderer recommends specializing unless you have a very clear concept and a very patient GM.

\textit{"I have tried to walk three roads at once. I ended up in a ditch, owing favors to a badger and singing hymns to a rock. Learn from my shame."}

\vspace{0.5cm}
\hrule
\vspace{0.2cm}
\begin{center}
\small \textit{"Choose your path. Pay its price. Tell the story."}
\end{center}

\chapter{Weaving the Threads -- Hybrid Magic Rules}

\epigraph{"I have seen a Runekeeper who learned to sing. I have met a Cantor who carried a dead spirit on her shoulder. I have watched a Witch touch a Psion's mind and both of them weep for the same lost memory. The roads cross. The question is not whether they can, but what you are willing to pay for the crossing."}{\textit{--- The Gray Wanderer}}

\section*{Why Hybrid Magic?}

The paths presented in this grimoire are not sealed vaults. A Cantor may envy the Invoker's flexibility. A Runekeeper may wish for the Witch's subtle touch. A Psion may hunger for the raw, dangerous flexibility of free casting.

Hybrid characters are possible, but they come with a cost: \textbf{divided attention}, \textbf{resource conversion inefficiency}, and \textbf{double jeopardy} (tracking two consequence tracks). These rules are not for every player. They are for those who want to tell a story of a character torn between two ways of power -- and who are willing to accept the friction.

\section*{General Hybrid Principles}

\begin{enumerate}
    \item \textbf{One Primary, One Secondary.} Choose which path is your character's main focus. The secondary path costs more XP to advance and its resources convert at a penalty.
    \item \textbf{Separate Tracks.} You track Obligation (Runekeeper/Invoker), Corruption (Cantor), Leash (Summoner), Promise Clocks (Witch), Mental Strain (Psion), and regular Fatigue separately. They do not merge.
    \item \textbf{Conversion Costs.} If a hybrid ability allows you to convert one resource to another, the rate is always \textbf{2:1} or worse. You cannot generate resources from nothing.
    \item \textbf{Hybrid Talents.} Specific talents (8--12 XP) unlock synergy between two paths. These are listed below.
\end{enumerate}

\section*{Resource Conversion Table}

Use these rates when a hybrid effect or GM ruling requires converting between resources. Conversions are always \textbf{one way} and cost an action unless noted.

% Replace the "Resource Exchange Rates" table with this:
\small
\begin{longtable}{p{2.5cm} p{2.5cm} p{1.5cm} p{5cm}}
\caption{Resource Exchange Rates (Hybrid Magic)}\\
\toprule
\textbf{From} & \textbf{To} & \textbf{Rate} & \textbf{Notes} \\
\midrule
Obligation (1 seg) & Fatigue & 2:1 & Meditation, once/session \\
Fatigue (2) & Obligation (1 seg) & 2:1 & Desperate bargain \\
Corruption (1 seg) & Mental Strain (1) & 1:1 & Cannot exceed cap \\
Mental Strain (2) & Corruption (1 seg) & 2:1 & Burnout taints soul \\
Leash (2 ticks) & Fatigue (1) & 2:1 & Spirit drains you \\
Promise Clock (1 seg) & Obligation (1 seg) & 1:1 & Hag bargain \\
Backlash (1 SB) & Mental Strain (1) & 1:1 & Crossfeed, unstable \\
\bottomrule
\end{longtable}
\normalsize

\noindent \textit{GM Note:} These conversions are intentionally inefficient. Hybrid characters should feel the strain of divided loyalties. The story is in the cost.

\section*{Hybrid Talents}

Each talent below costs \textbf{8 XP} unless noted. Prerequisites: at least one talent from each parent path, and experience using both in play (minimum 2 sessions with both active).

\subsection*{Runekeeper + Cantor -- The Hymn-Bound Codex}

\textit{"I sing the Rites. My Codex hums along."}

\textbf{Effect:} You may inscribe a Song (Low Rite only) into your Codex as if it were a Rite. Singing it requires a Performance + Lore roll (instead of Spirit + Lore) and costs the normal Obligation. The Song gains the [RESONANT] tag (it can trigger your Corruption clock if you Push it). In addition, when you use your Thiasos to aid a Performance roll, your familiar's assist dice count double for the purpose of group assists (max still +3).

\textbf{Downside:} Each time you invoke a Rite from your Codex immediately after using a Cantor's Inspire Chorus in the same scene, mark 1 additional Fatigue (the two powers resonate painfully).

\subsection*{Invoker + Witch -- The Crossroads Debtor}

\textit{"I make bargains with everyone. Patrons. Spirits. Hags. My own reflection."}

\textbf{Effect:} You may treat any Symbol you carry as a valid \textbf{threshold} for Witch rituals (doorway, altar, token). When you Crack the Seal on a Symbol, you may instead convert the Symbol's compromised state into a \textbf{Promise Clock [4]} (the debt becomes social/spiritual rather than physical damage). The Promise Clock ticks each time you use that Symbol without first making an offering (1 SB worth of wine, bread, or a whispered truth). If the Promise Clock fills, the Symbol shatters permanently.

\textbf{Downside:} You cannot clear Obligation and Promise Clocks in the same downtime scene. You must choose which master to serve first.

\subsection*{Summoner + Psion -- The Fractured Vessel}

\textit{"The spirit in my head is not a patron. It is a second self. And it is hungry."}

\textbf{Effect:} When you summon a spirit, you may choose to bind it using Mental Strain instead of a Leash. The spirit's Cap becomes your Psionics rating (min 1, max 3), and its Leash capacity is \texttt{Mental Strain capacity + Spirit}. Instead of ticking Leash, the spirit ticks your Mental Strain when it acts against its nature, takes harm, or crosses a ward. If Mental Strain overflows, the spirit immediately assumes control of your body for one exchange (GM narrates) before departing.

\textbf{Downside:} You cannot use any psionic Art while a spirit is bound -- your mind is too crowded. The exception is Mind Shield, which you may use at +1 Mental Strain cost to keep the spirit from influencing your thoughts.

\subsection*{Free Caster + Witch -- The Unwritten Covenant}

\textit{"I cast without Rites. I bind without Symbols. The world finds a way to charge me anyway."}

\textbf{Effect:} When you build a spell with TAGS, you may declare it a \textbf{threshold working}. The spell gains +1 Effect if cast at a proper threshold (doorway, crossroads, graveside, etc.) but also creates a \textbf{Promise Clock [2]} that ticks immediately after the spell resolves. On a miss or partial, the Clock ticks twice. When the Promise Clock fills, the world demands a price -- a lost memory, a broken friendship, a door that will not open for you. You may clear Promise Clock segments by making offerings (salt, bread, a truthful confession) before casting.

\textbf{Downside:} Any spell that includes the [DANGEROUS] tag (Leap, Fold, Gate, Create, Summon, Transmute, Animate) automatically creates a Promise Clock [4] regardless of success.

\subsection*{Cantor + Psion -- The Resonant Mind}

\textit{"I hear everyone's thoughts as a melody. The harmony is beautiful. The dissonance is agony."}

\textbf{Effect:} When you use Telepathy or Empathic Manipulation, you may add your Performance skill to the roll (instead of or in addition to, max +3 dice). You also may use your Corruption Clock as a \textbf{psychic battery}: once per session, you may advance your Corruption Clock by 1 segment to immediately recover 2 Mental Strain.

\textbf{Downside:} When your Corruption Clock fills, you do not get a Patron's boon -- instead, you have a \textbf{Mental Break}. Roll on the Psionic Complication table (d12) and apply the result immediately. Your Corruption Clock resets to 0, but you cannot use Cantor or Psionic abilities for the rest of the scene.

\subsection*{Runekeeper + Summoner -- The Stone Bond}

\textit{"My Thiasos is dead. The spirit in the Codex is not my familiar. It is my partner."}

\textbf{Effect:} You may bind a summoned spirit directly into your Codex (instead of a Thiasos). The spirit gains the Codex as an anchor; it does not depart when its Leash fills (instead, the Leash resets to 0 and you mark 2 Fatigue). The spirit can assist you in invoking Rites from the Codex, granting +1 die to the invocation. You may have only one such bound spirit at a time.

\textbf{Downside:} If the Codex is lost or damaged, the spirit is also lost -- and it will be \textit{angry}. The GM gains a 3 SB bank to spend on the spirit's return.

\subsection*{Witch + Psion -- The Silent Witness}

\textit{"I do not need to speak to be heard. I do not need to touch to feel. I am the threshold between minds."}

\textbf{Effect:} You may use your Mental Strain track as a substitute for \textbf{Soft Witness} (Witch talent). When you use Telepathy on a target who has accepted hospitality from you (a cup of tea, a night's shelter), the target's resistance DV is reduced by 1. Additionally, you may convert a Promise Clock segment into 1 Mental Strain recovery once per session (as a free action, but must be narrated as a small act of contrition).

\textbf{Downside:} Whenever you suffer Mental Strain overflow, tick a Promise Clock [4] (the universe notices you are a loose thread). This Promise Clock cannot be cleared except by a full downtime scene of meditation and apology.

\section*{Hybrid Build Examples}

\subsection*{The Inquisitor Psychic (Thepyrgos Witch-Hunter + Psion)}
\textit{"I hunt rogue mages. I also read their minds before they lie to me."}

Paths: Invoker (Inquisitor Prime) + Psion (Telepathy, Mind Shield)
Key Talents: Crossroads Debtor (Invoker+Witch reskin as Inquisitor+Psion -- GM approval), Mind Shield, Psychic Vampire
Playstyle: Use Telepathy to detect lies and track suspects. Use Mind Shield to protect against enemy psions. Use Invoker Rites (Cleansing Light, Marked Prey) to bind supernatural threats. Trade Symbol compromise for Promise Clocks -- then clear those clocks by testifying truthfully in court.

\subsection*The Spirit Cantor (Lethai Mask-Dancer + Summoner)}
\textit{"I wear the mask of the ancestor. The ancestor sings through me."}

Paths: Cantor + Summoner (ancestral spirits)
Key Talents: The Hymn-Bound Codex (reskinned as "The Mask-Bound Spirit"), Spirit Link
Playstyle: Summon a wolf ancestor or selkie. Use the Cantor's Inspire Chorus to buff allies while the spirit attacks. When the spirit acts against its nature, you tick your Corruption Clock instead of Leash -- you are becoming the mask.

\subsection*The Psionic Witch (Silkstrand Hedge-Witch + Clairvoyance)}
\textit{"I see the threads. I pluck them. Sometimes the thread plucks back."}

Paths: Witch + Psion (Clairvoyance, Telepathy)
Key Talents: The Silent Witness, Weaver of Remembering
Playstyle: Use Clairvoyance to scout ahead and find hidden thresholds. Use Witch rites to bind promises and bless hearths. Convert Promise Clock segments into Mental Strain to keep scrying. Never accept hospitality from someone you might later need to read -- the debt complicates both.

\section*{GM Guidance for Hybrids}

\begin{itemize}
    \item \textbf{Require narrative justification.} A character does not wake up one day as a Runekeeper/Psion. They found an ancient text. They survived a psychic storm. They made a bargain with a spirit that rewired their brain. Play that story.
    \item \textbf{Be generous with complications.} Hybrids are magnets for interesting trouble. A Patron might resent the character's other powers. A Promise Clock might tick when they use telepathy without consent. A summoned spirit might be jealous of a Cantor's audience.
    \item \textbf{Watch the conversion economy.} If a player finds a conversion rate that feels broken (e.g., converting free Backlash SB into cheap Mental Strain), veto it. The rates above are playtested but not perfect. Trust your gut.
    \item \textbf{Offer hybrid-specific story hooks.} The Gray Wanderer encountered a Psion whose mind had been colonized by a Carrion-King. A Cantor whose voice awakened a dormant spirit. A Summoner who became a Witch's Hollowed. These are campaign seeds, not just builds.
\end{itemize}

\section*{A Final Word on Power}

Hybrid characters are not more powerful than single-path characters. They are more \textit{versatile}, but their versatility comes with double the resource tracking and double the consequence tracks. In a typical session, a pure Runekeeper will invoke three Rites and manage one Obligation clock. A hybrid Runekeeper/Psion will invoke two Rites, use two psionic Arts, and manage Obligation, Mental Strain, and probably a Promise Clock. The mental load is real.

Play hybrid if you want the story of a character pulled between two philosophies. Play hybrid if you love bookkeeping and resource conversion puzzles. Play hybrid if you want your GM to have endless ammunition for complications.

Do not play hybrid if you want to be more powerful than everyone else. The Gray Wanderer tried that. He woke up in a ditch, owing favors to a badger.

\vspace{0.5cm}
\hrule
\vspace{0.2cm}
\begin{center}
\small \textit{"Weave carefully. The threads are watching."}
\end{center}

\end{document}
